% ==============================================================
% Chapter: Labor and Capital
% 

\epigraph{As means are constantly being found for the maintenance of labor on cheaper and more wretched food, the minimum of wages is constantly sinking.}{\textit{Karl Marx}, ``On the Question of Free Trade'' (1847)}

\epigraph{Three-hour shifts or a fifteen-hour week may put off the problem for a great while. For three hours a day is quite enough to satisfy the old Adam in most of us!}{\textit{John Maynard Keynes}, ``Economic Possibilities for our Grandchildren'' (1930)}

==============================================================
\label{chap:labor-capital}

This chapter applies the generalized means framework to the interaction between labor and capital---the central factors in production and growth theory. We begin with a two-period model that clarifies how the elasticity of substitution between factors governs the response to technological change and factor price movements. We then extend to nested structures with multiple capital types, dynamic accumulation, and stochastic environments. The chapter culminates with the Real Business Cycle model, showing how our CES toolkit illuminates the canonical workhorse of modern macroeconomics, and then moves beyond to incorporate factor heterogeneity.

% ==============================================================
\section{The Two-Period Problem with Labor and Capital}
\label{sec:lc-two-period}
% ==============================================================

We begin by studying factor choice in a two-period setting. A firm combines labor and capital to produce output, and the elasticity of substitution between these factors determines how the firm responds to changes in factor prices and technology. This section establishes the core framework that will be extended to dynamics and uncertainty in subsequent sections.

% --------------------------------------------------------------
\subsection{From Varieties to Factors}
\label{subsec:varieties-to-factors}
% --------------------------------------------------------------

In the canonical CES problem, a consumer chose quantities $\bx = (x_1, \ldots, x_n)$ of different goods subject to a budget constraint. The relative price $p_i/p_j$ governed the trade-off between goods $i$ and $j$, and the elasticity of substitution $\varepsilon$ determined how sensitive the allocation was to price changes.

Now consider a firm that combines two factors---labor $L$ and capital $K$---to produce output $Y$. The fundamental trade-off is no longer between consumption goods but between hiring workers and renting machines. The wage $w$ and the rental rate $r$ play the role of prices, and output plays the role of utility.

\paragraph{The Production Function as a Generalized Mean.}
If the firm's technology exhibits constant elasticity of substitution, output is a generalized mean of factor inputs:
\begin{equation}
Y = \F(K, L; \bom, \sigma) = A \cdot \M(K, L; \bom, \sigma),
\label{eq:ces-production}
\end{equation}
where $A > 0$ is total factor productivity (TFP), $\bom = (\omega_K, \omega_L)$ with $\omega_K + \omega_L = 1$ are factor weights, and $\sigma > 0$ is the elasticity of substitution between capital and labor.

In canonical form, the production function is:
\begin{equation}
Y = A \left( \omega_K^{1-\rho} K^\rho + \omega_L^{1-\rho} L^\rho \right)^{1/\rho}, \qquad \rho = 1 - \frac{1}{\sigma}.
\label{eq:ces-canonical}
\end{equation}

\paragraph{The Cost Minimization Problem.}
The firm's problem is the dual of the consumer's: minimize the cost of producing a target output level $\bar{Y}$:
\begin{equation}
\min_{K, L} \left\{ rK + wL \right\} \quad \text{s.t.} \quad \F(K, L) \geq \bar{Y}.
\label{eq:cost-min}
\end{equation}

This is exactly the structure of the canonical problem with the roles reversed: output is the constraint, and cost is the objective. The solutions from \cref{sec:problem-solution} apply directly.

% --------------------------------------------------------------
\subsection{The Canonical Structure and Primal-Dual Symmetry}
\label{subsec:canonical-production}
% --------------------------------------------------------------

The connection between production and the canonical CES problem runs deeper than analogy. We now make this precise, showing that cost minimization and utility maximization are dual problems with symmetric solutions.

\paragraph{The Primal Problem: Utility Maximization.}
Recall the consumer's problem from \cref{sec:problem-solution}:
\begin{equation}
\max_{\bx} \; \M(\bx; \bom, \rho) \quad \text{s.t.} \quad \sum_i p_i x_i = B.
\label{eq:primal-utility}
\end{equation}
The solution yields optimal quantities $x_i^* = \omega_i (p_i / \Pstar)^{-\varepsilon} \cdot (B / \Pstar)$, where the price index is:
\begin{equation}
\Pstar = \M_{1-\varepsilon}(\bp; \bom) = \left( \sum_i \omega_i \, p_i^{1-\varepsilon} \right)^{1/(1-\varepsilon)}.
\label{eq:price-index-recall}
\end{equation}
The indirect utility is $\M^* = B / \Pstar$: the budget divided by the price index.

\paragraph{The Dual Problem: Cost Minimization.}
The firm's cost minimization problem is:
\begin{equation}
\min_{K, L} \; \left\{ rK + wL \right\} \quad \text{s.t.} \quad A \cdot \M(K, L; \bom, \rho) \geq \bar{Y}.
\label{eq:dual-cost}
\end{equation}
This is the \emph{dual} of the primal: we minimize cost (the ``budget'') subject to achieving a target utility (output).

\paragraph{The Unit Cost Index as a Price Index.}
Solving the cost minimization problem yields the cost function:
\begin{equation}
C(\bar{Y}, r, w) = \Cstar \cdot \frac{\bar{Y}}{A},
\label{eq:cost-function}
\end{equation}
where the \emph{unit cost index} is:
\begin{equation}
\Cstar = \M_{1-\sigma}(r, w; \bom) = \left( \omega_K \, r^{1-\sigma} + \omega_L \, w^{1-\sigma} \right)^{1/(1-\sigma)}.
\label{eq:unit-cost-as-mean}
\end{equation}

\begin{calloutimportant}[The Unit Cost is a Generalized Mean of Factor Prices]
The unit cost index $\Cstar$ is itself a generalized mean---specifically, a mean of factor prices $(r, w)$ with weights $(\omega_K, \omega_L)$ and power parameter $1 - \sigma$. This is the \emph{same} formula as the ideal price index in the consumption problem, with factor prices replacing goods prices and $\sigma$ replacing $\varepsilon$.

The symmetry is complete:
\begin{center}
\renewcommand{\arraystretch}{1.3}
\begin{tabular}{@{}lcc@{}}
\toprule
& \textbf{Consumption} & \textbf{Production} \\
\midrule
Quantities & Goods $x_i$ & Factors $K, L$ \\
Prices & Goods prices $p_i$ & Factor prices $r, w$ \\
Objective & Maximize utility $\M(\bx)$ & Minimize cost $rK + wL$ \\
Constraint & Budget $\sum p_i x_i = B$ & Output $\M(K,L) \geq \bar{Y}/A$ \\
Index & Price index $\Pstar$ & Unit cost $\Cstar$ \\
Value & $\M^* = B/\Pstar$ & $C^* = \Cstar \cdot \bar{Y}/A$ \\
\bottomrule
\end{tabular}
\end{center}
\end{calloutimportant}

\paragraph{Factor Shares as Induced Weights.}
In the consumption problem, the expenditure share of good $i$ is $\wexp_i = p_i x_i / B$. At the optimum, this equals the induced weight:
\begin{equation}
\wexp_i = \omega_i \left( \frac{p_i}{\Pstar} \right)^{1-\varepsilon}.
\label{eq:exp-share-recall}
\end{equation}

The production analog is the \emph{cost share}---the fraction of total cost paid to each factor:
\begin{equation}
s_K = \frac{rK}{rK + wL} = \omega_K \left( \frac{r}{\Cstar} \right)^{1-\sigma}, \qquad s_L = \omega_L \left( \frac{w}{\Cstar} \right)^{1-\sigma}.
\label{eq:cost-share-induced}
\end{equation}

These are exactly the induced weights of the generalized mean $\Cstar$ evaluated at factor prices $(r, w)$. The cost shares emerge from the same mathematical structure as expenditure shares.

\paragraph{The Transformation Perspective.}
Recall from \cref{sec:transformations} that the canonical and classical representations are related by a weight transformation. In the production context, we can write the unit cost in classical form:
\begin{equation}
\Cstar = \left( \tilde{\omega}_K \, r^{1-\sigma} + \tilde{\omega}_L \, w^{1-\sigma} \right)^{1/(1-\sigma)},
\label{eq:unit-cost-classical}
\end{equation}
where $\tilde{\omega}_K + \tilde{\omega}_L = 1$ are the transformed weights. The cost shares are then:
\begin{equation}
s_K = \frac{\tilde{\omega}_K \, r^{1-\sigma}}{\tilde{\omega}_K \, r^{1-\sigma} + \tilde{\omega}_L \, w^{1-\sigma}}, \qquad s_L = \frac{\tilde{\omega}_L \, w^{1-\sigma}}{\tilde{\omega}_K \, r^{1-\sigma} + \tilde{\omega}_L \, w^{1-\sigma}}.
\label{eq:shares-classical}
\end{equation}

This is the ``inner sum divided by total'' formula for induced weights. The cost shares are the normalized contributions to the inner sum of the unit cost aggregator.

\begin{calloutnote}[Why the Canonical Form?]
The canonical representation $\M(K, L; \bom, \rho)$ with weights $\omega_i^{1-\rho}$ inside the sum has a key advantage: when factors are paid their marginal products, the cost shares $s_i$ equal the induced weights of the production function evaluated at the optimum. This connects factor payments directly to the structure of technology.

In contrast, the classical form requires tracking transformed weights that depend on the power parameter. The canonical form keeps the primitive weights $\omega_i$ interpretable across different values of $\sigma$.
\end{calloutnote}

% --------------------------------------------------------------
\subsection{Complementary Factors}
\label{subsec:complementary-factors}
% --------------------------------------------------------------

When $\sigma < 1$, capital and labor are \emph{complements}: an increase in one factor raises the marginal product of the other. The limiting case $\sigma \to 0$ yields the Leontief technology, where factors must be used in fixed proportions.

\paragraph{The Cobb-Douglas Benchmark.}
The knife-edge case $\sigma = 1$ (equivalently, $\rho \to 0$) gives the Cobb-Douglas production function:
\begin{equation}
Y = A \, K^{\omega_K} L^{\omega_L}.
\label{eq:cobb-douglas}
\end{equation}
This specification has remarkable properties. Factor shares are constant regardless of factor prices:
\begin{equation}
\text{Capital share:} \quad \frac{rK}{Y} = \omega_K, \qquad \text{Labor share:} \quad \frac{wL}{Y} = \omega_L.
\label{eq:cd-shares}
\end{equation}

The Cobb-Douglas case serves as the dividing line between complements ($\sigma < 1$) and substitutes ($\sigma > 1$).

\paragraph{Factor Demands.}
Applying the solution formulas from the canonical problem, the cost-minimizing factor demands are:
\begin{equation}
K^* = \omega_K \left( \frac{r}{\Cstar} \right)^{-\sigma} \cdot \frac{\bar{Y}/A}{\Cstar}, \qquad L^* = \omega_L \left( \frac{w}{\Cstar} \right)^{-\sigma} \cdot \frac{\bar{Y}/A}{\Cstar},
\label{eq:factor-demands}
\end{equation}
where $\Cstar$ is the unit cost index (the cost of producing one unit of output):
\begin{equation}
\Cstar = \left( \omega_K \, r^{1-\sigma} + \omega_L \, w^{1-\sigma} \right)^{1/(1-\sigma)}.
\label{eq:unit-cost}
\end{equation}

\paragraph{The Cost Share Formulas.}
The share of total cost paid to each factor is:
\begin{equation}
s_K \equiv \frac{rK}{rK + wL} = \omega_K \left( \frac{r}{\Cstar} \right)^{1-\sigma}, \qquad s_L = \omega_L \left( \frac{w}{\Cstar} \right)^{1-\sigma}.
\label{eq:cost-shares}
\end{equation}

When $\sigma < 1$, a rise in the relative price of capital $(r/w)$ \emph{increases} the capital share. The intuition is that with complementary factors, the firm cannot easily substitute away from the more expensive input, so its expenditure share rises.

\begin{calloutimportant}[Complements: Price and Share Move Together]
When factors are complements ($\sigma < 1$), the factor whose relative price rises sees its cost share \emph{increase}. This is the opposite of what happens with substitutes. The firm is ``stuck'' using the expensive factor because it cannot produce without it.
\end{calloutimportant}

% --------------------------------------------------------------
\subsection{Substitutable Factors}
\label{subsec:substitutable-factors}
% --------------------------------------------------------------

When $\sigma > 1$, capital and labor are \emph{substitutes}: the firm can more easily replace one factor with the other in response to price changes. This case has attracted renewed attention due to concerns about automation and the declining labor share.

\paragraph{High Elasticity of Substitution.}
With $\sigma > 1$, a rise in the relative price of labor $(w/r)$ leads to a \emph{fall} in the labor share. The firm substitutes toward capital so aggressively that labor's share of costs declines despite the higher wage.

\paragraph{The Labor Share and Automation.}
Consider the response of the labor share to a fall in the rental rate of capital (perhaps due to technological progress in producing capital goods):
\begin{equation}
\frac{d \log s_L}{d \log r} = (1 - \sigma)(1 - s_L).
\label{eq:labor-share-response}
\end{equation}

When $\sigma > 1$, this expression is negative: cheaper capital reduces the labor share. This mechanism is central to theories of automation and the ``rise of the robots.''

\begin{calloutnote}[The Automation Debate]
The empirical debate centers on the magnitude of $\sigma$. If $\sigma$ is close to or below one, cheaper capital goods need not threaten labor's share. If $\sigma$ substantially exceeds one, technological progress that reduces the cost of capital can significantly compress labor income as a fraction of output.

Cross-country and time-series evidence suggests $\sigma$ may vary by sector and time period, with some estimates above one for the aggregate economy in recent decades.
\end{calloutnote}

\paragraph{Comparative Statics Summary.}
The following table summarizes how factor shares respond to a rise in the relative price of capital $(r/w)$:

\begin{center}
\renewcommand{\arraystretch}{1.3}
\begin{tabular}{@{}lcc@{}}
\toprule
& $\sigma < 1$ (Complements) & $\sigma > 1$ (Substitutes) \\
\midrule
Capital share $s_K$ & Increases & Decreases \\
Labor share $s_L$ & Decreases & Increases \\
Capital-labor ratio $K/L$ & Decreases & Decreases \\
\bottomrule
\end{tabular}
\end{center}

The capital-labor ratio always falls when capital becomes relatively more expensive---this is the substitution effect. But the \emph{share} response depends on whether the quantity adjustment is large enough to offset the price change.

% --------------------------------------------------------------
\subsection{Types of Technical Change}
\label{subsec:technical-change}
% --------------------------------------------------------------

Technical progress can take different forms depending on which factor it augments. The classification matters for long-run growth and for the evolution of factor shares.

\paragraph{The Augmented Production Function.}
We generalize \cref{eq:ces-canonical} to allow factor-specific productivity:
\begin{equation}
Y = \left( \omega_K^{1-\rho} (A_K K)^\rho + \omega_L^{1-\rho} (A_L L)^\rho \right)^{1/\rho},
\label{eq:augmented-ces}
\end{equation}
where $A_K$ is capital-augmenting technology and $A_L$ is labor-augmenting technology. The type of technical change is classified by which productivity term grows over time.

\paragraph{Harrod-Neutral (Labor-Augmenting) Technical Change.}
Technical progress is \emph{Harrod-neutral} if it augments labor: $A_L$ grows while $A_K$ is constant. The production function becomes:
\begin{equation}
Y = \F(K, A_L L).
\label{eq:harrod-neutral}
\end{equation}

This is the only form of technical change consistent with balanced growth when $\sigma \neq 1$. Along a balanced growth path with Harrod-neutral progress, the capital-output ratio $K/Y$ is constant, the interest rate is constant, and output per worker grows at the rate of labor-augmenting technical change.

\paragraph{Solow-Neutral (Capital-Augmenting) Technical Change.}
Technical progress is \emph{Solow-neutral} if it augments capital: $A_K$ grows while $A_L$ is constant. The production function becomes:
\begin{equation}
Y = \F(A_K K, L).
\label{eq:solow-neutral}
\end{equation}

With $\sigma \neq 1$, Solow-neutral progress is \emph{not} consistent with balanced growth. The capital share trends over time, and either capital or labor eventually dominates production.

\paragraph{Hicks-Neutral Technical Change.}
Technical progress is \emph{Hicks-neutral} if it scales output without changing the marginal rate of technical substitution:
\begin{equation}
Y = A \cdot \F(K, L),
\label{eq:hicks-neutral}
\end{equation}
where $A$ grows over time. This is equivalent to setting $A_K = A_L = A^{1/\rho}$ in the augmented formulation.

With Cobb-Douglas technology ($\sigma = 1$), all three types of neutral technical change are equivalent. With $\sigma \neq 1$, they have distinct implications.

\begin{calloutimportant}[Why Harrod-Neutrality Matters]
The requirement for balanced growth with $\sigma \neq 1$ restricts attention to labor-augmenting technical change. This is a strong assumption, and its empirical validity is debated. The restriction disappears in the Cobb-Douglas case, which partly explains that specification's enduring popularity in growth models.
\end{calloutimportant}

\paragraph{Effects on Factor Shares.}
Consider Harrod-neutral progress (rising $A_L$). The effective labor supply is $\tilde{L} = A_L L$, so the labor share becomes:
\begin{equation}
s_L = \omega_L \left( \frac{w/A_L}{\Cstar} \right)^{1-\sigma}.
\label{eq:share-harrod}
\end{equation}

If $\sigma < 1$, labor-augmenting progress \emph{raises} the labor share: effective labor becomes more abundant relative to capital, but with complementary factors, labor's share rises. If $\sigma > 1$, the opposite occurs.

The following table summarizes the effect of factor-augmenting technical progress on factor shares:

\begin{center}
\renewcommand{\arraystretch}{1.3}
\begin{tabular}{@{}lcc@{}}
\toprule
Type of Progress & Effect on $s_L$ if $\sigma < 1$ & Effect on $s_L$ if $\sigma > 1$ \\
\midrule
Labor-augmenting ($\uparrow A_L$) & $\uparrow$ & $\downarrow$ \\
Capital-augmenting ($\uparrow A_K$) & $\downarrow$ & $\uparrow$ \\
\bottomrule
\end{tabular}
\end{center}

% --------------------------------------------------------------
\subsection{The Dual Problem: Profit Maximization}
\label{subsec:dual-problem}
% --------------------------------------------------------------

Rather than minimizing cost for a given output, the firm may maximize profit taking factor prices as given. With constant returns to scale, the two problems are closely related.

\paragraph{The Profit Function.}
A competitive firm taking output price $P$, wage $w$, and rental rate $r$ as given solves:
\begin{equation}
\max_{K, L} \left\{ P \cdot \F(K, L) - rK - wL \right\}.
\label{eq:profit-max}
\end{equation}

With constant returns to scale, the first-order conditions equate factor prices to marginal products:
\begin{equation}
r = P \cdot \frac{\partial \F}{\partial K}, \qquad w = P \cdot \frac{\partial \F}{\partial L}.
\label{eq:foc-profit}
\end{equation}

\paragraph{Marginal Products in CES.}
From the canonical CES production function \cref{eq:ces-canonical}:
\begin{align}
\frac{\partial Y}{\partial K} &= A \cdot \omega_K^{1-\rho} \left( \frac{Y}{AK} \right)^{1-\rho} = \omega_K^{1-\rho} A^\rho \left( \frac{Y}{K} \right)^{1-\rho}, \label{eq:mpk} \\[6pt]
\frac{\partial Y}{\partial L} &= A \cdot \omega_L^{1-\rho} \left( \frac{Y}{AL} \right)^{1-\rho} = \omega_L^{1-\rho} A^\rho \left( \frac{Y}{L} \right)^{1-\rho}. \label{eq:mpl}
\end{align}

The marginal product of each factor depends on the output-to-factor ratio raised to the power $1-\rho = 1/\sigma$. When $\sigma < 1$ (complements), marginal products are highly sensitive to factor ratios. When $\sigma > 1$ (substitutes), marginal products are less sensitive.

\paragraph{Zero Profit with Constant Returns.}
With constant returns to scale, Euler's theorem implies:
\begin{equation}
Y = \frac{\partial Y}{\partial K} \cdot K + \frac{\partial Y}{\partial L} \cdot L.
\label{eq:euler}
\end{equation}

Combined with the first-order conditions \cref{eq:foc-profit}, this yields zero economic profit:
\begin{equation}
P \cdot Y = rK + wL \quad \Longrightarrow \quad \text{Profit} = 0.
\label{eq:zero-profit}
\end{equation}

This is the production-side analog of the budget constraint holding with equality in the consumer problem.

% --------------------------------------------------------------
\subsection{Summary}
\label{subsec:two-period-lc-summary}
% --------------------------------------------------------------

\begin{table}[H]
\centering
\caption{The CES Production Framework}
\label{tab:ces-production-summary}
\renewcommand{\arraystretch}{1.6}
\begin{tabular}{@{}ll@{}}
\toprule
\textbf{Object} & \textbf{Formula} \\
\midrule
Production function & $Y = A \left( \omega_K^{1-\rho} K^\rho + \omega_L^{1-\rho} L^\rho \right)^{1/\rho}$ \\[6pt]
Elasticity of substitution & $\sigma = 1/(1-\rho)$ \\[6pt]
Unit cost index & $\Cstar = \left( \omega_K \, r^{1-\sigma} + \omega_L \, w^{1-\sigma} \right)^{1/(1-\sigma)}$ \\[6pt]
Capital share & $s_K = \omega_K \left( r / \Cstar \right)^{1-\sigma}$ \\[6pt]
Labor share & $s_L = \omega_L \left( w / \Cstar \right)^{1-\sigma}$ \\[6pt]
\bottomrule
\end{tabular}
\end{table}

\begin{callouttip}[Key Insight]
The two-period labor-capital problem is a CES optimization problem in the factor space:
\begin{itemize}[nosep]
\item Factors $\{K, L\}$ play the role of goods $\{1, 2\}$
\item Factor prices $(r, w)$ are the relative prices
\item The elasticity of substitution $\sigma$ governs factor substitution
\item The weights $(\omega_K, \omega_L)$ capture factor importance
\end{itemize}
All the machinery of generalized means---cost indices, factor shares, comparative statics---applies directly. The key economic question is the magnitude of $\sigma$: complements ($\sigma < 1$) versus substitutes ($\sigma > 1$) yield qualitatively different responses to price and technology changes.
\end{callouttip}


% ==============================================================
\section{Two Capital Stocks: Nested CES}
\label{sec:nested-capital}
% ==============================================================

Many production processes involve multiple types of capital that interact differently with labor. Structures and equipment, for example, may have distinct substitution patterns with workers. This section extends the framework to nested CES aggregators, where an outer layer combines labor with a capital composite, and an inner layer aggregates different capital types. The nested structure allows us to capture rich substitution patterns while maintaining tractability.

% --------------------------------------------------------------
\subsection{The Nested Structure}
\label{subsec:nested-structure}
% --------------------------------------------------------------

Consider a firm that uses labor $L$ and two types of capital: equipment $K_E$ and structures $K_S$. Rather than treating all three as symmetric inputs, we allow for a hierarchical aggregation.

\paragraph{The Two-Level Aggregator.}
The production function has a nested structure:
\begin{equation}
Y = \F^{\text{outer}}\bigl(L, \, \F^{\text{inner}}(K_E, K_S)\bigr),
\label{eq:nested-production}
\end{equation}
where the inner aggregator combines the two capital types into a capital composite $K$, and the outer aggregator combines labor with this composite.

\paragraph{Inner Layer: Capital Aggregation.}
The capital composite is a CES aggregate of equipment and structures:
\begin{equation}
K = \M(K_E, K_S; \bom^K, \sigma_K) = \left( \omega_E^{1-\rho_K} K_E^{\rho_K} + \omega_S^{1-\rho_K} K_S^{\rho_K} \right)^{1/\rho_K},
\label{eq:capital-composite}
\end{equation}
where $\sigma_K = 1/(1-\rho_K)$ is the elasticity of substitution between capital types, and $\omega_E + \omega_S = 1$.

\paragraph{Outer Layer: Labor-Capital Aggregation.}
Output combines labor and the capital composite:
\begin{equation}
Y = A \cdot \M(K, L; \bom, \sigma) = A \left( \omega_K^{1-\rho} K^{\rho} + \omega_L^{1-\rho} L^{\rho} \right)^{1/\rho},
\label{eq:outer-aggregator}
\end{equation}
where $\sigma = 1/(1-\rho)$ is the elasticity of substitution between labor and the capital composite, and $\omega_K + \omega_L = 1$.

\begin{calloutimportant}[Complementary Outer, Substitutable Inner]
A leading specification sets $\sigma < 1$ (labor and capital are complements) and $\sigma_K > 1$ (equipment and structures are substitutes). This captures the intuition that:
\begin{itemize}[nosep]
\item Workers and machines must work together---you cannot produce with capital alone
\item But within the capital stock, firms can flexibly substitute between different types of machinery and buildings
\end{itemize}
This asymmetry has important implications for how technological progress in different sectors affects factor shares.
\end{calloutimportant}

\paragraph{The Expanded Form.}
Substituting the capital composite into the outer layer yields the full production function:
\begin{equation}
Y = A \left( \omega_K^{1-\rho} \left[ \omega_E^{1-\rho_K} K_E^{\rho_K} + \omega_S^{1-\rho_K} K_S^{\rho_K} \right]^{\rho/\rho_K} + \omega_L^{1-\rho} L^{\rho} \right)^{1/\rho}.
\label{eq:nested-expanded}
\end{equation}

This is a three-factor production function with a specific substitution structure encoded in the nesting.

% --------------------------------------------------------------
\subsection{Nested Price Indices as Nested Means}
\label{subsec:nested-price-indices}
% --------------------------------------------------------------

The nested production structure induces a corresponding nested structure in price indices. This connection---between quantity aggregation and price aggregation---is a powerful feature of the CES framework.

\paragraph{The Duality Principle.}
Recall from \cref{subsec:canonical-production} that the unit cost index is a generalized mean of factor prices:
\begin{equation}
\Cstar = \M_{1-\sigma}(\text{factor prices}; \bom).
\label{eq:duality-principle}
\end{equation}
With nested production, this principle applies recursively: each level of quantity aggregation has a corresponding level of price aggregation.

\paragraph{Inner Price Index: The Cost of Capital.}
The capital composite $K = \M(K_E, K_S; \bom^K, \sigma_K)$ aggregates equipment and structures. Its ``price''---the unit cost of the capital composite---is the dual mean:
\begin{equation}
C_K^* = \M_{1-\sigma_K}(r_E, r_S; \bom^K) = \left( \omega_E \, r_E^{1-\sigma_K} + \omega_S \, r_S^{1-\sigma_K} \right)^{1/(1-\sigma_K)}.
\label{eq:inner-price-index}
\end{equation}

This is a generalized mean of rental rates with the \emph{same weights} $(\omega_E, \omega_S)$ as the quantity aggregator, but with power parameter $1 - \sigma_K$ instead of $\rho_K = 1 - 1/\sigma_K$.

\paragraph{Outer Price Index: Total Unit Cost.}
The outer aggregator $Y = A \cdot \M(K, L; \bom, \sigma)$ combines the capital composite with labor. Its dual is the total unit cost:
\begin{equation}
\Cstar = \M_{1-\sigma}(C_K^*, w; \bom) = \left( \omega_K \, (C_K^*)^{1-\sigma} + \omega_L \, w^{1-\sigma} \right)^{1/(1-\sigma)}.
\label{eq:outer-price-index}
\end{equation}

The total unit cost is a generalized mean of the capital price index $C_K^*$ and the wage $w$---a mean of means.

\begin{calloutimportant}[Nested Means All the Way Down]
The nested structure creates a hierarchy of generalized means:
\begin{align*}
\text{Quantities:} \quad & K = \M_{\sigma_K}(K_E, K_S) \quad \longrightarrow \quad Y = A \cdot \M_\sigma(K, L) \\[4pt]
\text{Prices:} \quad & C_K^* = \M_{1-\sigma_K}(r_E, r_S) \quad \longrightarrow \quad \Cstar = \M_{1-\sigma}(C_K^*, w)
\end{align*}
The quantity aggregators use power $\rho = 1 - 1/\sigma$; the price aggregators use power $1 - \sigma$. This duality ensures that the cost function has the multiplicative form $C = \Cstar \cdot Y/A$.
\end{calloutimportant}

\paragraph{Shares as Induced Weights at Each Level.}
The induced weight interpretation extends to each level of the nest. At the inner level:
\begin{equation}
\frac{s_E}{s_E + s_S} = \omega_E \left( \frac{r_E}{C_K^*} \right)^{1-\sigma_K}, \qquad \frac{s_S}{s_E + s_S} = \omega_S \left( \frac{r_S}{C_K^*} \right)^{1-\sigma_K}.
\label{eq:inner-shares}
\end{equation}
These are the induced weights of the inner price index $C_K^*$.

At the outer level:
\begin{equation}
s_K = \omega_K \left( \frac{C_K^*}{\Cstar} \right)^{1-\sigma}, \qquad s_L = \omega_L \left( \frac{w}{\Cstar} \right)^{1-\sigma}.
\label{eq:outer-shares}
\end{equation}
These are the induced weights of the outer price index $\Cstar$.

The total equipment share is the product: $s_E = s_K \cdot (s_E/(s_E + s_S))$---the share of cost going to capital, times the share of capital cost going to equipment.

\paragraph{The Recursive Decomposition.}
This structure connects to the recursive decomposition of generalized means from \cref{sec:recursive-decomposition}. Define the ``effective prices'' as:
\begin{equation}
\tilde{r}_E \equiv \frac{r_E}{C_K^*}, \qquad \tilde{r}_S \equiv \frac{r_S}{C_K^*}, \qquad \tilde{C}_K \equiv \frac{C_K^*}{\Cstar}, \qquad \tilde{w} \equiv \frac{w}{\Cstar}.
\label{eq:effective-prices}
\end{equation}

Then the cost shares have the uniform representation:
\begin{equation}
s_i = \omega_i \, \tilde{p}_i^{1-\sigma_i},
\label{eq:uniform-shares}
\end{equation}
where $\tilde{p}_i$ is the effective price (price relative to the relevant index) and $\sigma_i$ is the elasticity at that level of the nest.

\begin{calloutnote}[Connection to Discrete Choice]
The nested CES structure is isomorphic to the nested logit model in discrete choice. In the logit framework, choice probabilities take the form $\pi_i \propto \exp(v_i / \lambda)$, where $v_i$ is the utility of option $i$ and $\lambda$ is a scale parameter. The nested logit allows different scale parameters at different levels of the choice hierarchy.

The CES analog has shares $s_i \propto \omega_i \, p_i^{1-\sigma}$. The elasticity $\sigma$ plays the role of the inverse scale parameter $1/\lambda$. Nested CES allows different elasticities at different levels, just as nested logit allows different scales. This connection is developed further in \cref{sec:discrete-choice}.
\end{calloutnote}

% --------------------------------------------------------------
\subsection{Factor Prices and Cost Minimization}
\label{subsec:nested-cost-min}
% --------------------------------------------------------------

Let $r_E$ and $r_S$ denote the rental rates of equipment and structures, and let $w$ denote the wage. The firm minimizes the cost of producing output $\bar{Y}$.

\paragraph{The Inner Problem.}
For a given level of the capital composite $K$, the firm minimizes the cost of capital:
\begin{equation}
\min_{K_E, K_S} \left\{ r_E K_E + r_S K_S \right\} \quad \text{s.t.} \quad \M(K_E, K_S; \bom^K, \sigma_K) \geq K.
\label{eq:inner-cost-min}
\end{equation}

Applying the standard CES solution, the cost-minimizing demands are:
\begin{equation}
K_E^* = \omega_E \left( \frac{r_E}{C_K^*} \right)^{-\sigma_K} \cdot \frac{K}{C_K^*}, \qquad K_S^* = \omega_S \left( \frac{r_S}{C_K^*} \right)^{-\sigma_K} \cdot \frac{K}{C_K^*},
\label{eq:inner-demands}
\end{equation}
where the unit cost of the capital composite is:
\begin{equation}
C_K^* = \left( \omega_E \, r_E^{1-\sigma_K} + \omega_S \, r_S^{1-\sigma_K} \right)^{1/(1-\sigma_K)}.
\label{eq:capital-unit-cost}
\end{equation}

\paragraph{The Outer Problem.}
Taking $C_K^*$ as the ``price'' of the capital composite, the firm then solves:
\begin{equation}
\min_{K, L} \left\{ C_K^* \cdot K + w \cdot L \right\} \quad \text{s.t.} \quad A \cdot \M(K, L; \bom, \sigma) \geq \bar{Y}.
\label{eq:outer-cost-min}
\end{equation}

The solution gives:
\begin{equation}
K^* = \omega_K \left( \frac{C_K^*}{\Cstar} \right)^{-\sigma} \cdot \frac{\bar{Y}/A}{\Cstar}, \qquad L^* = \omega_L \left( \frac{w}{\Cstar} \right)^{-\sigma} \cdot \frac{\bar{Y}/A}{\Cstar},
\label{eq:outer-demands}
\end{equation}
where the overall unit cost index is:
\begin{equation}
\Cstar = \left( \omega_K \, (C_K^*)^{1-\sigma} + \omega_L \, w^{1-\sigma} \right)^{1/(1-\sigma)}.
\label{eq:total-unit-cost}
\end{equation}

\paragraph{Cost Shares.}
The share of total cost paid to each factor is:
\begin{align}
s_L &= \omega_L \left( \frac{w}{\Cstar} \right)^{1-\sigma}, \label{eq:labor-share-nested} \\[6pt]
s_K &= \omega_K \left( \frac{C_K^*}{\Cstar} \right)^{1-\sigma} = 1 - s_L, \label{eq:capital-share-nested}
\end{align}
and within the capital share:
\begin{align}
s_E &= s_K \cdot \omega_E \left( \frac{r_E}{C_K^*} \right)^{1-\sigma_K}, \label{eq:equipment-share} \\[6pt]
s_S &= s_K \cdot \omega_S \left( \frac{r_S}{C_K^*} \right)^{1-\sigma_K}. \label{eq:structures-share}
\end{align}

The nested structure implies that the equipment share of total cost, $s_E$, depends on both the outer elasticity $\sigma$ (which governs how the capital share responds to the capital price index) and the inner elasticity $\sigma_K$ (which governs the equipment-structures mix).

% --------------------------------------------------------------
\subsection{One-Off Technological Progress}
\label{subsec:tech-progress}
% --------------------------------------------------------------

We now examine how a one-time improvement in technology affects factor demands and shares. The analysis reveals how the nesting structure shapes the transmission of shocks.

\paragraph{Progress in the Equipment Sector.}
Suppose equipment-augmenting technology improves: each unit of equipment becomes more productive, which we model as a fall in the effective rental rate $r_E$. What happens to factor shares?

The direct effect operates through the inner layer. With $\sigma_K > 1$ (substitutes), a fall in $r_E$ leads to:
\begin{itemize}[nosep]
\item A rise in equipment use relative to structures: $K_E/K_S$ increases
\item A fall in the equipment share of capital costs: $s_E/(s_E + s_S)$ decreases
\item A fall in the capital price index $C_K^*$
\end{itemize}

The indirect effect operates through the outer layer. The fall in $C_K^*$ makes the capital composite cheaper relative to labor. With $\sigma < 1$ (complements):
\begin{itemize}[nosep]
\item The capital-labor ratio $K/L$ rises, but not by much
\item The capital share $s_K$ \emph{falls} (cheaper input, higher share---complements)
\item The labor share $s_L$ \emph{rises}
\end{itemize}

\begin{calloutnote}[Equipment-Biased Technical Change]
With the complementary-outer, substitutable-inner specification, equipment-biased technical change \emph{raises} the labor share. The intuition: cheaper equipment reduces the cost of the capital composite, but with labor-capital complementarity, this benefits workers through higher labor demand. The substitutability within capital means the cost savings are large, amplifying the effect.
\end{calloutnote}

\paragraph{Progress in the Structures Sector.}
Now suppose structures-augmenting technology improves (a fall in $r_S$). The logic is symmetric:
\begin{itemize}[nosep]
\item Inner layer: structures expand relative to equipment; structures share of capital falls; $C_K^*$ falls
\item Outer layer: labor share rises (same mechanism as above)
\end{itemize}

The key difference is the magnitude. If equipment is a larger share of the capital composite ($\omega_E > \omega_S$ or $s_E > s_S$), then progress in equipment has a larger effect on $C_K^*$ and hence a larger effect on the labor share.

\paragraph{Comparative Statics Summary.}
Define the elasticity of the labor share with respect to the equipment rental rate:
\begin{equation}
\eta_{s_L, r_E} \equiv \frac{\partial \log s_L}{\partial \log r_E}.
\label{eq:labor-share-elasticity}
\end{equation}

With $\sigma < 1$ and $\sigma_K > 1$, this elasticity is negative: cheaper equipment raises the labor share. The magnitude depends on:
\begin{enumerate}[nosep]
\item The outer elasticity $\sigma$: closer to zero means stronger complementarity, larger effect
\item The inner elasticity $\sigma_K$: larger means more substitution within capital, larger cost reduction
\item The equipment weight $\omega_E$ and share $s_E$: larger means equipment matters more for $C_K^*$
\end{enumerate}

\begin{proposition}[Labor Share Response to Equipment Progress]
\label{prop:labor-share-equipment}
With $\sigma < 1$ and $\sigma_K > 1$, the elasticity of the labor share with respect to the equipment rental rate is:
\begin{equation}
\eta_{s_L, r_E} = -(1 - \sigma) \cdot s_K \cdot \frac{s_E}{s_K} = -(1 - \sigma) \cdot s_E.
\label{eq:eta-formula}
\end{equation}
A one percent fall in the equipment rental rate raises the labor share by $(1-\sigma) \cdot s_E$ percent.
\end{proposition}

% --------------------------------------------------------------
\subsection{Asymmetric Nesting: Labor-Equipment Complementarity}
\label{subsec:asymmetric-nesting}
% --------------------------------------------------------------

An alternative specification places labor and equipment in the inner nest, with structures in the outer layer. This captures ``capital-skill complementarity'' when we interpret equipment as skill-intensive capital.

\paragraph{The Alternative Structure.}
\begin{equation}
Y = A \cdot \M\Bigl(K_S, \, \M(K_E, L; \bom^{EL}, \sigma_{EL}); \bom, \sigma_S\Bigr),
\label{eq:alt-nested}
\end{equation}
where the inner layer combines equipment and labor with elasticity $\sigma_{EL}$, and the outer layer combines this composite with structures.

\paragraph{Capital-Skill Complementarity.}
If $\sigma_{EL} < 1$, equipment and labor are complements. Progress in equipment (cheaper $r_E$) then has different effects:
\begin{itemize}[nosep]
\item Inner layer: equipment expands; the equipment-labor composite becomes cheaper
\item With complements, the labor share \emph{within} the inner composite rises
\item The overall effect on the labor share depends on how the inner composite interacts with structures
\end{itemize}

This specification is relevant for understanding skill-biased technical change, where computerization (equipment) complements skilled labor but substitutes for unskilled labor.

% --------------------------------------------------------------
\subsection{Summary}
\label{subsec:nested-summary}
% --------------------------------------------------------------

\begin{table}[H]
\centering
\caption{Nested CES Production}
\label{tab:nested-summary}
\renewcommand{\arraystretch}{1.6}
\begin{tabular}{@{}ll@{}}
\toprule
\textbf{Object} & \textbf{Formula} \\
\midrule
Capital composite & $K = \left( \omega_E^{1-\rho_K} K_E^{\rho_K} + \omega_S^{1-\rho_K} K_S^{\rho_K} \right)^{1/\rho_K}$ \\[6pt]
Capital price index & $C_K^* = \left( \omega_E \, r_E^{1-\sigma_K} + \omega_S \, r_S^{1-\sigma_K} \right)^{1/(1-\sigma_K)}$ \\[6pt]
Output & $Y = A \left( \omega_K^{1-\rho} K^{\rho} + \omega_L^{1-\rho} L^{\rho} \right)^{1/\rho}$ \\[6pt]
Total unit cost & $\Cstar = \left( \omega_K \, (C_K^*)^{1-\sigma} + \omega_L \, w^{1-\sigma} \right)^{1/(1-\sigma)}$ \\[6pt]
Labor share & $s_L = \omega_L \left( w / \Cstar \right)^{1-\sigma}$ \\[6pt]
\bottomrule
\end{tabular}
\end{table}

\begin{callouttip}[Key Insight]
The nested CES structure allows different substitution patterns at different levels of aggregation:
\begin{itemize}[nosep]
\item Outer elasticity $\sigma$: governs labor-capital substitution
\item Inner elasticity $\sigma_K$: governs substitution across capital types
\end{itemize}
With complements at the outer layer ($\sigma < 1$) and substitutes at the inner layer ($\sigma_K > 1$), technological progress that reduces the cost of any capital type \emph{raises} the labor share. The nested structure amplifies this effect by allowing large within-capital substitution.

This framework provides a foundation for analyzing how sector-specific technological change---such as the falling price of computing equipment---affects the distribution of income between labor and capital.
\end{callouttip}


% ==============================================================
\section{Dynamics with Labor and Capital}
\label{sec:lc-dynamics}
% ==============================================================

We now extend the labor-capital framework to a dynamic setting where capital is accumulated over time. The firm or household chooses how much to invest each period, trading off current consumption against future productive capacity. The elasticity of substitution between labor and capital continues to play a central role, now shaping not only static factor demands but also the dynamics of capital accumulation and the transition to steady state.

% --------------------------------------------------------------
\subsection{The Infinite Horizon Problem}
\label{subsec:infinite-horizon}
% --------------------------------------------------------------

Consider a representative household that owns capital, supplies labor, and makes consumption-saving decisions over an infinite horizon. Output is produced using a CES technology, and the household chooses consumption $\{c_t\}_{t=0}^\infty$ and capital $\{k_{t+1}\}_{t=0}^\infty$ to maximize lifetime utility.

\paragraph{Preferences.}
The household has time-separable preferences with discount factor $\beta \in (0,1)$ and period utility $u(c) = \frac{c^{1-1/\theta} - 1}{1 - 1/\theta}$, where $\theta > 0$ is the intertemporal elasticity of substitution (IES):
\begin{equation}
\V_0 = \sum_{t=0}^{\infty} \beta^t u(c_t) = \sum_{t=0}^{\infty} \beta^t \frac{c_t^{1-1/\theta} - 1}{1 - 1/\theta}.
\label{eq:lifetime-utility}
\end{equation}

\paragraph{Technology.}
Output is produced using capital $k_t$ and labor $\ell_t$ according to the CES production function:
\begin{equation}
y_t = A \left( \omega_K^{1-\rho} k_t^\rho + \omega_L^{1-\rho} (B_t \ell_t)^\rho \right)^{1/\rho},
\label{eq:ces-production-dynamic}
\end{equation}
where $A$ is total factor productivity, $B_t$ is labor-augmenting technology (growing at rate $g$: $B_{t+1} = (1+g) B_t$), and $\sigma = 1/(1-\rho)$ is the elasticity of substitution.

\paragraph{Capital Accumulation.}
Capital depreciates at rate $\delta \in (0,1)$ and is augmented by investment $i_t$:
\begin{equation}
k_{t+1} = (1-\delta) k_t + i_t.
\label{eq:capital-accumulation}
\end{equation}

\paragraph{Resource Constraint.}
Output is divided between consumption and investment:
\begin{equation}
c_t + i_t = y_t.
\label{eq:resource-constraint}
\end{equation}

Substituting the capital accumulation equation:
\begin{equation}
c_t + k_{t+1} = y_t + (1-\delta) k_t.
\label{eq:resource-consolidated}
\end{equation}

\paragraph{Labor Supply.}
For simplicity, we assume labor is supplied inelastically at $\ell_t = 1$ for all $t$. This allows us to focus on the capital accumulation margin.

% --------------------------------------------------------------
\subsection{The Euler Equation}
\label{subsec:euler-equation}
% --------------------------------------------------------------

The household's problem is to maximize \cref{eq:lifetime-utility} subject to \cref{eq:resource-consolidated} and the production technology \cref{eq:ces-production-dynamic}. The key optimality condition is the Euler equation.

\paragraph{Derivation.}
The Lagrangian for the household's problem is:
\begin{equation}
\mathcal{L} = \sum_{t=0}^{\infty} \beta^t \left[ u(c_t) + \lambda_t \left( y_t + (1-\delta) k_t - c_t - k_{t+1} \right) \right].
\label{eq:lagrangian-dynamic}
\end{equation}

The first-order conditions with respect to $c_t$ and $k_{t+1}$ are:
\begin{align}
c_t: \quad & u'(c_t) = \lambda_t, \label{eq:foc-c} \\[4pt]
k_{t+1}: \quad & \lambda_t = \beta \lambda_{t+1} \left( \frac{\partial y_{t+1}}{\partial k_{t+1}} + 1 - \delta \right). \label{eq:foc-k}
\end{align}

Combining these yields the Euler equation:
\begin{equation}
\boxed{u'(c_t) = \beta \, u'(c_{t+1}) \left( \frac{\partial y_{t+1}}{\partial k_{t+1}} + 1 - \delta \right).}
\label{eq:euler-equation}
\end{equation}

\paragraph{The Marginal Product of Capital.}
From the CES production function \cref{eq:ces-production-dynamic}:
\begin{equation}
\frac{\partial y_t}{\partial k_t} = A^\rho \, \omega_K^{1-\rho} \left( \frac{y_t}{k_t} \right)^{1-\rho} = \omega_K^{1-\rho} A^\rho \left( \frac{y_t}{k_t} \right)^{1/\sigma}.
\label{eq:mpk-dynamic}
\end{equation}

The marginal product depends on the output-capital ratio raised to the power $1/\sigma$. When $\sigma < 1$ (complements), the marginal product is highly sensitive to the capital-output ratio. When $\sigma > 1$ (substitutes), it is less sensitive.

\paragraph{The Euler Equation in Growth Rates.}
With CRRA utility, $u'(c) = c^{-1/\theta}$, the Euler equation becomes:
\begin{equation}
\left( \frac{c_{t+1}}{c_t} \right)^{1/\theta} = \beta \left( \frac{\partial y_{t+1}}{\partial k_{t+1}} + 1 - \delta \right).
\label{eq:euler-growth}
\end{equation}

Defining the gross return on capital as $R_{t+1} \equiv \frac{\partial y_{t+1}}{\partial k_{t+1}} + 1 - \delta$:
\begin{equation}
\frac{c_{t+1}}{c_t} = \left( \beta R_{t+1} \right)^\theta.
\label{eq:consumption-growth}
\end{equation}

This is the dynamic analog of the consumption ratio formula from the two-period problem. Consumption growth is increasing in the return to capital, with elasticity equal to the IES $\theta$.

% --------------------------------------------------------------
\subsection{The Intertemporal CES Structure}
\label{subsec:intertemporal-ces}
% --------------------------------------------------------------

The dynamic problem inherits the CES structure developed in earlier chapters. We now make this connection explicit, showing how the Euler equation and value function relate to the canonical framework.

\paragraph{Time as Another Dimension of Choice.}
In \cref{sec:lc-two-period}, we showed that the firm's choice between capital and labor is a CES optimization problem. The consumer's choice between consumption today and tomorrow is \emph{also} a CES problem---time replaces factors as the dimension of choice.

Consider the two-period version of \cref{eq:lifetime-utility}:
\begin{equation}
\V = u(c_0) + \beta \, u(c_1) = \frac{c_0^{1-1/\theta}}{1-1/\theta} + \beta \frac{c_1^{1-1/\theta}}{1-1/\theta}.
\label{eq:two-period-utility}
\end{equation}

After the monotone transformation to generalized mean form (as in \cref{sec:cs-two-period}), this becomes:
\begin{equation}
\V \sim \M(c_0, c_1; \bom^\beta, \theta),
\label{eq:utility-as-mean}
\end{equation}
where the weights $\bom^\beta = (\omega_0, \omega_1)$ are derived from the discount factor $\beta$.

\paragraph{The Intertemporal Price Index.}
The budget constraint $c_0 + R^{-1} c_1 = W$ has prices $(1, R^{-1})$, where $R^{-1}$ is the price of future consumption in terms of current consumption. The ideal price index is:
\begin{equation}
\Pstar_{\text{time}} = \M_{1-\theta}(1, R^{-1}; \bom^\beta) = \left( \omega_0 + \omega_1 \, R^{-(1-\theta)} \right)^{1/(1-\theta)}.
\label{eq:intertemporal-price-index}
\end{equation}

This is a generalized mean of intertemporal prices with the same structure as the unit cost index in production.

\paragraph{The Euler Equation as a Relative Price Condition.}
The first-order condition for the two-period problem equates the marginal rate of substitution to the price ratio:
\begin{equation}
\frac{u'(c_0)}{u'(c_1)} = \frac{1}{R^{-1}} = R.
\label{eq:mrs-price}
\end{equation}

With CRRA utility, $u'(c) = c^{-1/\theta}$, this becomes:
\begin{equation}
\left( \frac{c_0}{c_1} \right)^{-1/\theta} = R \quad \Longrightarrow \quad \frac{c_1}{c_0} = R^\theta.
\label{eq:consumption-ratio-static}
\end{equation}

Including the discount factor, the full condition is:
\begin{equation}
\frac{u'(c_0)}{\beta \, u'(c_1)} = R \quad \Longrightarrow \quad \frac{c_1}{c_0} = (\beta R)^\theta.
\label{eq:euler-from-ratio}
\end{equation}

This is exactly the Euler equation \cref{eq:consumption-growth}. The dynamic first-order condition is the CES ratio condition applied intertemporally.

\begin{calloutimportant}[Three CES Problems, One Structure]
The chapter has developed three distinct CES optimization problems:
\begin{center}
\renewcommand{\arraystretch}{1.3}
\begin{tabular}{@{}lccc@{}}
\toprule
& \textbf{Factors} & \textbf{Capital Types} & \textbf{Time} \\
\midrule
Quantities & $K, L$ & $K_E, K_S$ & $c_0, c_1, \ldots$ \\
Prices & $r, w$ & $r_E, r_S$ & $1, R^{-1}, R^{-2}, \ldots$ \\
Elasticity & $\sigma$ & $\sigma_K$ & $\theta$ \\
Price index & $\Cstar$ & $C_K^*$ & $\Pstar_{\text{time}}$ \\
Shares & $s_K, s_L$ & $s_E, s_S$ & $\wexp_0, \wexp_1, \ldots$ \\
\bottomrule
\end{tabular}
\end{center}
All three share the same mathematical structure: a generalized mean objective, a linear constraint, and solutions characterized by the ratio condition and induced weights.
\end{calloutimportant}

\paragraph{The Value Function and Homogeneity.}
The value function for the infinite-horizon problem inherits the homogeneity of the CES structure. With CRRA utility and constant returns in production, the value function takes the form:
\begin{equation}
V(k) = \phi \cdot \frac{k^{1-1/\theta}}{1 - 1/\theta},
\label{eq:value-function-form}
\end{equation}
where $\phi$ is a constant that depends on parameters $(\beta, \delta, A, \omega_K, \omega_L, \sigma, \theta)$.

This homogeneity implies that the policy functions are linear in capital:
\begin{equation}
c^*(k) = (1 - s) \cdot y(k), \qquad i^*(k) = s \cdot y(k),
\label{eq:policy-linear}
\end{equation}
where $s$ is the (endogenous) saving rate. The saving rate is constant along the balanced growth path, varying only during transitions.

\paragraph{Connection to the Canonical Problem.}
The infinite-horizon problem can be viewed as a limit of finite-horizon CES problems. With $T$ periods, lifetime utility is:
\begin{equation}
\V_T = \sum_{t=0}^{T} \beta^t u(c_t) \sim \M(c_0, c_1, \ldots, c_T; \bom^\beta_T, \theta),
\label{eq:finite-horizon}
\end{equation}
where the weights $\bom^\beta_T$ are derived from the discount factors $(\beta^0, \beta^1, \ldots, \beta^T)$.

As $T \to \infty$, the finite-horizon problem converges to the infinite-horizon problem. The recursive structure (Bellman equation) emerges as the limiting form of the nested CES aggregation---each period's continuation value is itself a generalized mean of future consumption.

\begin{calloutnote}[The Epstein-Zin Connection]
With time-separable utility, the IES $\theta$ also governs attitudes toward risk (through the coefficient of relative risk aversion $\gamma = 1/\theta$). Epstein-Zin preferences break this link by using separate parameters for intertemporal substitution and risk aversion.

In the Epstein-Zin framework, the value function is a CES aggregate of current consumption and the certainty equivalent of future utility:
\[
V_t = \M(c_t, \CE(V_{t+1}); \bom^\beta, \theta),
\]
where $\CE(\cdot)$ is a generalized mean over states with power parameter related to risk aversion. This ``CES over CES'' structure is developed in \cref{sec:epstein-zin}.
\end{calloutnote}

% --------------------------------------------------------------
\subsection{Steady State}
\label{subsec:steady-state}
% --------------------------------------------------------------

A steady state (or balanced growth path) is an allocation where all variables grow at constant rates. With labor-augmenting technical change at rate $g$, the balanced growth path has output, capital, and consumption all growing at rate $g$.

\paragraph{Stationarity in Efficiency Units.}
Define variables in efficiency units by dividing by $B_t$:
\begin{equation}
\tilde{k}_t \equiv \frac{k_t}{B_t}, \qquad \tilde{c}_t \equiv \frac{c_t}{B_t}, \qquad \tilde{y}_t \equiv \frac{y_t}{B_t}.
\label{eq:efficiency-units}
\end{equation}

In efficiency units, the production function becomes:
\begin{equation}
\tilde{y}_t = A \left( \omega_K^{1-\rho} \tilde{k}_t^\rho + \omega_L^{1-\rho} \right)^{1/\rho},
\label{eq:production-efficiency}
\end{equation}
which is time-invariant when $\tilde{k}_t$ is constant.

\paragraph{The Steady-State Capital Stock.}
In steady state, consumption grows at rate $g$, so $c_{t+1}/c_t = 1 + g$. The Euler equation \cref{eq:consumption-growth} implies:
\begin{equation}
(1+g)^{1/\theta} = \beta R^* \quad \Longrightarrow \quad R^* = \frac{(1+g)^{1/\theta}}{\beta}.
\label{eq:steady-state-return}
\end{equation}

The steady-state return $R^*$ is pinned down by preferences and the growth rate, independent of the production technology. This is the modified golden rule.

\paragraph{Solving for Steady-State Capital.}
The return on capital is $R^* = MPK^* + 1 - \delta$, where $MPK^*$ is the steady-state marginal product. Using \cref{eq:mpk-dynamic}:
\begin{equation}
MPK^* = R^* - 1 + \delta = \omega_K^{1-\rho} A^\rho \left( \frac{\tilde{y}^*}{\tilde{k}^*} \right)^{1/\sigma}.
\label{eq:mpk-steady-state}
\end{equation}

This equation, together with the production function \cref{eq:production-efficiency}, determines the steady-state capital-output ratio $\tilde{k}^*/\tilde{y}^*$.

\begin{calloutnote}[The Role of $\sigma$ in Steady State]
The elasticity of substitution $\sigma$ affects the \emph{level} of steady-state capital but not the steady-state \emph{return} (which is pinned down by the Euler equation). With $\sigma < 1$ (complements), a given capital-output ratio implies a higher marginal product of capital; the steady-state capital stock is lower. With $\sigma > 1$ (substitutes), the marginal product is less sensitive to capital intensity; the steady-state capital stock is higher.
\end{calloutnote}

\paragraph{The Steady-State Capital Share.}
The capital share in steady state is:
\begin{equation}
s_K^* = \frac{MPK^* \cdot \tilde{k}^*}{\tilde{y}^*} = \omega_K^{1-\rho} A^\rho \left( \frac{\tilde{y}^*}{\tilde{k}^*} \right)^{1/\sigma - 1} = \omega_K \left( \frac{MPK^*}{\omega_K^{1-\rho} A^\rho} \right)^{1 - \sigma}.
\label{eq:capital-share-steady}
\end{equation}

With Cobb-Douglas ($\sigma = 1$), the capital share equals $\omega_K$ regardless of the level of capital or technology. With $\sigma \neq 1$, the capital share depends on factor prices and technology.

% --------------------------------------------------------------
\subsection{Transition Dynamics}
\label{subsec:transition-dynamics}
% --------------------------------------------------------------

Away from steady state, the economy transitions toward the balanced growth path. The speed and nature of this transition depend critically on the elasticity of substitution.

\paragraph{The Phase Diagram.}
The dynamics of the economy can be characterized by a system in $(\tilde{k}_t, \tilde{c}_t)$. The resource constraint in efficiency units is:
\begin{equation}
\tilde{c}_t + (1+g) \tilde{k}_{t+1} = \tilde{y}_t + (1-\delta) \tilde{k}_t.
\label{eq:resource-efficiency}
\end{equation}

The Euler equation in efficiency units becomes:
\begin{equation}
\frac{\tilde{c}_{t+1}}{\tilde{c}_t} = \frac{1}{1+g} \left( \beta R_{t+1} \right)^\theta.
\label{eq:euler-efficiency}
\end{equation}

\paragraph{The $\dot{k} = 0$ Locus.}
Capital is constant ($\tilde{k}_{t+1} = \tilde{k}_t$) when:
\begin{equation}
\tilde{c} = \tilde{y}(\tilde{k}) - (g + \delta) \tilde{k}.
\label{eq:k-dot-zero}
\end{equation}

This is the familiar condition that consumption equals output minus the investment required to maintain capital per efficiency unit.

\paragraph{The $\dot{c} = 0$ Locus.}
Consumption is constant ($\tilde{c}_{t+1} = \tilde{c}_t$) when the return equals its steady-state value:
\begin{equation}
R(\tilde{k}) = MPK(\tilde{k}) + 1 - \delta = R^* = \frac{(1+g)^{1/\theta}}{\beta}.
\label{eq:c-dot-zero}
\end{equation}

This pins down a unique value of $\tilde{k}$ (the steady-state capital stock).

\paragraph{Dynamics with Complements vs. Substitutes.}
The speed of convergence to steady state depends on how rapidly the marginal product of capital changes as capital accumulates.

With $\sigma < 1$ (complements): the marginal product is highly sensitive to capital. Starting below steady state, the high return induces rapid capital accumulation, but the return falls quickly as capital rises. Convergence is relatively fast.

With $\sigma > 1$ (substitutes): the marginal product is less sensitive to capital. The return remains elevated for longer as capital accumulates, sustaining high investment. However, the ``target'' is harder to reach precisely because diminishing returns set in slowly. Convergence can be slower.

\begin{calloutimportant}[Convergence Speed and $\sigma$]
The local speed of convergence near steady state is governed by the eigenvalues of the linearized system. A key determinant is the elasticity of the marginal product with respect to capital:
\begin{equation}
\frac{\partial \log MPK}{\partial \log k} = \frac{1}{\sigma} - 1 = -\frac{1-\sigma}{\sigma}.
\label{eq:mpk-elasticity}
\end{equation}
When $\sigma < 1$, this elasticity is large in magnitude (strong diminishing returns), implying faster convergence. When $\sigma > 1$, the elasticity is smaller, implying slower convergence.
\end{calloutimportant}

\paragraph{Numerical Illustration.}
Consider two economies that differ only in $\sigma$, starting from the same initial capital stock below steady state.

\begin{center}
\renewcommand{\arraystretch}{1.3}
\begin{tabular}{@{}lcc@{}}
\toprule
& $\sigma = 0.5$ (Complements) & $\sigma = 2$ (Substitutes) \\
\midrule
Initial MPK & High & Moderately high \\
Initial investment rate & Very high & High \\
MPK decline as $k$ rises & Rapid & Gradual \\
Half-life to steady state & Short & Long \\
\bottomrule
\end{tabular}
\end{center}

The complementary economy exhibits more dramatic swings: high initial returns drive a burst of investment, but returns fall quickly, moderating the investment boom. The substitutable economy has a more gradual transition.

% --------------------------------------------------------------
\subsection{Capital Accumulation with Two Capital Types}
\label{subsec:two-capital-dynamics}
% --------------------------------------------------------------

The dynamic framework extends naturally to the nested CES case with multiple capital types. Let $k_E$ denote equipment and $k_S$ denote structures, each with its own accumulation equation:
\begin{align}
k_{E,t+1} &= (1-\delta_E) k_{E,t} + i_{E,t}, \label{eq:equipment-accumulation} \\[4pt]
k_{S,t+1} &= (1-\delta_S) k_{S,t} + i_{S,t}. \label{eq:structures-accumulation}
\end{align}

\paragraph{Two Euler Equations.}
The household now faces two investment margins. The Euler equations are:
\begin{align}
u'(c_t) &= \beta \, u'(c_{t+1}) \left( \frac{\partial y_{t+1}}{\partial k_{E,t+1}} + 1 - \delta_E \right), \label{eq:euler-equipment} \\[4pt]
u'(c_t) &= \beta \, u'(c_{t+1}) \left( \frac{\partial y_{t+1}}{\partial k_{S,t+1}} + 1 - \delta_S \right). \label{eq:euler-structures}
\end{align}

\paragraph{Equalization of Returns.}
In equilibrium, the two Euler equations imply that the returns on equipment and structures are equalized:
\begin{equation}
\frac{\partial y}{\partial k_E} + 1 - \delta_E = \frac{\partial y}{\partial k_S} + 1 - \delta_S.
\label{eq:return-equalization}
\end{equation}

If depreciation rates differ ($\delta_E > \delta_S$ for equipment), the marginal product of equipment must exceed that of structures to compensate.

\paragraph{Transition with Sector-Specific Shocks.}
Suppose there is a one-time improvement in equipment technology (a permanent fall in the effective price of equipment). The economy transitions to a new steady state with:
\begin{itemize}[nosep]
\item Higher equipment-to-structures ratio
\item Higher capital composite relative to labor
\item Potentially different factor shares (depending on elasticities)
\end{itemize}

The transition dynamics involve both capital types adjusting, with equipment rising faster initially due to the direct shock.

% --------------------------------------------------------------
\subsection{Summary}
\label{subsec:dynamics-summary}
% --------------------------------------------------------------

\begin{table}[H]
\centering
\caption{Dynamic Labor-Capital Framework}
\label{tab:dynamics-summary}
\renewcommand{\arraystretch}{1.6}
\begin{tabular}{@{}ll@{}}
\toprule
\textbf{Object} & \textbf{Formula} \\
\midrule
Euler equation & $u'(c_t) = \beta \, u'(c_{t+1}) \left( MPK_{t+1} + 1 - \delta \right)$ \\[6pt]
Consumption growth & $c_{t+1}/c_t = \left( \beta R_{t+1} \right)^\theta$ \\[6pt]
Marginal product (CES) & $MPK = \omega_K^{1-\rho} A^\rho \left( y/k \right)^{1/\sigma}$ \\[6pt]
Steady-state return & $R^* = (1+g)^{1/\theta} / \beta$ \\[6pt]
MPK elasticity & $\partial \log MPK / \partial \log k = (1-\sigma)/\sigma$ \\[6pt]
\bottomrule
\end{tabular}
\end{table}

\begin{callouttip}[Key Insight]
The elasticity of substitution $\sigma$ governs the dynamics of capital accumulation:
\begin{itemize}[nosep]
\item With complements ($\sigma < 1$): strong diminishing returns, rapid convergence, volatile investment
\item With substitutes ($\sigma > 1$): weak diminishing returns, slow convergence, smooth investment
\end{itemize}
The steady-state return is pinned down by preferences and growth, independent of $\sigma$. But the steady-state capital stock, factor shares, and transition dynamics all depend critically on the degree of substitutability between labor and capital.

This framework provides the foundation for the Real Business Cycle model, where we add stochastic shocks and study fluctuations around the balanced growth path.
\end{callouttip}


% ==============================================================
\section{The Real Business Cycle Model}
\label{sec:rbc}
% ==============================================================

The Real Business Cycle (RBC) model is the canonical framework for studying macroeconomic fluctuations. It extends the neoclassical growth model by introducing stochastic productivity shocks, generating business cycle dynamics from the optimal responses of households and firms. This section develops the RBC model through the lens of CES aggregation, showing how the elasticity of substitution between labor and capital shapes the transmission of shocks and the behavior of factor shares over the cycle.

% --------------------------------------------------------------
\subsection{The Stochastic Growth Model}
\label{subsec:stochastic-growth}
% --------------------------------------------------------------

The RBC model adds uncertainty to the dynamic framework of \cref{sec:lc-dynamics}. Productivity is no longer deterministic but follows a stochastic process, and agents form expectations over future states.

\paragraph{Preferences.}
The representative household maximizes expected lifetime utility:
\begin{equation}
\V_0 = \E_0 \left[ \sum_{t=0}^{\infty} \beta^t u(c_t, \ell_t) \right],
\label{eq:rbc-utility}
\end{equation}
where $\E_0$ denotes the expectation conditional on information at date $0$. We now allow for endogenous labor supply, with period utility:
\begin{equation}
u(c, \ell) = \frac{c^{1-1/\theta}}{1 - 1/\theta} - \chi \frac{\ell^{1+1/\nu}}{1 + 1/\nu},
\label{eq:period-utility}
\end{equation}
where $\theta > 0$ is the IES, $\nu > 0$ is the Frisch elasticity of labor supply, and $\chi > 0$ scales the disutility of labor.

\paragraph{Technology.}
Output is produced using the CES production function:
\begin{equation}
y_t = A_t \left( \omega_K^{1-\rho} k_t^\rho + \omega_L^{1-\rho} \ell_t^\rho \right)^{1/\rho},
\label{eq:rbc-production}
\end{equation}
where $A_t$ is stochastic total factor productivity (TFP). For now, we consider Hicks-neutral shocks; we discuss factor-biased shocks below.

\paragraph{The Productivity Process.}
TFP follows an AR(1) process in logs:
\begin{equation}
\log A_t = \bar{a} + \rho_A (\log A_{t-1} - \bar{a}) + \varepsilon_t, \qquad \varepsilon_t \sim N(0, \sigma_A^2),
\label{eq:tfp-process}
\end{equation}
where $\bar{a}$ is the unconditional mean, $\rho_A \in (0,1)$ is the persistence, and $\sigma_A$ is the volatility of innovations.

\paragraph{Capital Accumulation and Resource Constraint.}
Capital evolves as before:
\begin{equation}
k_{t+1} = (1 - \delta) k_t + i_t, \qquad c_t + i_t = y_t.
\label{eq:rbc-constraints}
\end{equation}

\paragraph{The Household's Problem.}
The household chooses sequences $\{c_t, \ell_t, k_{t+1}\}$ to maximize \cref{eq:rbc-utility} subject to \cref{eq:rbc-constraints} and the production technology \cref{eq:rbc-production}, taking the stochastic process for $A_t$ as given.

% --------------------------------------------------------------
\subsection{Equilibrium Conditions}
\label{subsec:rbc-equilibrium}
% --------------------------------------------------------------

The first-order conditions characterize the equilibrium. We derive them and highlight how the CES structure shapes each condition.

\paragraph{The Stochastic Euler Equation.}
The intertemporal optimality condition becomes:
\begin{equation}
u_c(c_t, \ell_t) = \beta \, \E_t \left[ u_c(c_{t+1}, \ell_{t+1}) \cdot R_{t+1} \right],
\label{eq:stochastic-euler}
\end{equation}
where $R_{t+1} = \frac{\partial y_{t+1}}{\partial k_{t+1}} + 1 - \delta$ is the (stochastic) gross return on capital.

With the utility specification \cref{eq:period-utility}, $u_c = c^{-1/\theta}$, so:
\begin{equation}
c_t^{-1/\theta} = \beta \, \E_t \left[ c_{t+1}^{-1/\theta} \cdot R_{t+1} \right].
\label{eq:euler-crra}
\end{equation}

\paragraph{The Labor Supply Condition.}
The intratemporal optimality condition equates the marginal rate of substitution between consumption and leisure to the wage:
\begin{equation}
\frac{u_\ell(c_t, \ell_t)}{u_c(c_t, \ell_t)} = w_t,
\label{eq:labor-supply}
\end{equation}
where $w_t = \frac{\partial y_t}{\partial \ell_t}$ is the marginal product of labor.

With our utility specification:
\begin{equation}
\chi \ell_t^{1/\nu} \cdot c_t^{1/\theta} = w_t.
\label{eq:labor-supply-explicit}
\end{equation}

This can be rewritten as a labor supply curve:
\begin{equation}
\ell_t = \left( \frac{w_t}{\chi \, c_t^{1/\theta}} \right)^\nu.
\label{eq:labor-supply-curve}
\end{equation}

Labor supply increases with the wage (elasticity $\nu$) and decreases with consumption (wealth effect).

\paragraph{The Marginal Products in CES Form.}
From the production function \cref{eq:rbc-production}:
\begin{align}
R_t &= A_t^\rho \, \omega_K^{1-\rho} \left( \frac{y_t}{k_t} \right)^{1/\sigma} + 1 - \delta, \label{eq:return-ces} \\[6pt]
w_t &= A_t^\rho \, \omega_L^{1-\rho} \left( \frac{y_t}{\ell_t} \right)^{1/\sigma}. \label{eq:wage-ces}
\end{align}

The marginal products depend on output-to-factor ratios raised to the power $1/\sigma$. This is where the elasticity of substitution enters the equilibrium conditions.

\begin{calloutnote}[The Role of $\sigma$ in RBC Dynamics]
The elasticity $\sigma$ affects how marginal products respond to factor ratios:
\begin{itemize}[nosep]
\item With $\sigma < 1$ (complements): marginal products are highly sensitive to factor ratios. A productivity shock that raises output relative to capital causes a large increase in the return to capital.
\item With $\sigma > 1$ (substitutes): marginal products are less sensitive. The same shock causes a smaller increase in returns.
\end{itemize}
This shapes the volatility of investment, the persistence of output, and the comovement of factor prices over the cycle.
\end{calloutnote}

% --------------------------------------------------------------
\subsection{Productivity Shocks and Business Cycle Dynamics}
\label{subsec:rbc-dynamics}
% --------------------------------------------------------------

We now analyze how productivity shocks propagate through the economy. The key insight is that the CES structure governs the amplification and persistence of shocks.

\paragraph{Impact of a Positive TFP Shock.}
Consider an unexpected increase in $A_t$. The immediate effects are:

\begin{enumerate}[nosep]
\item \textbf{Output rises}: Higher TFP directly increases $y_t$ for given inputs.
\item \textbf{Marginal products rise}: Both $R_t$ and $w_t$ increase, as factor productivities are enhanced.
\item \textbf{Labor supply increases}: The higher wage induces more work (substitution effect dominates).
\item \textbf{Consumption rises}: Higher income and wealth raise consumption.
\item \textbf{Investment rises}: The higher return to capital makes saving attractive.
\end{enumerate}

\paragraph{The Amplification Mechanism.}
The labor supply response amplifies the initial shock. Higher TFP raises wages, which increases labor supply, which further raises output. The strength of this amplification depends on:
\begin{itemize}[nosep]
\item The Frisch elasticity $\nu$: higher $\nu$ means more responsive labor supply
\item The elasticity of substitution $\sigma$: affects how wages respond to labor input
\end{itemize}

With $\sigma < 1$ (complements), an increase in labor raises the marginal product of capital and lowers the marginal product of labor (diminishing returns are strong). This dampens the wage response and limits amplification.

With $\sigma > 1$ (substitutes), diminishing returns are weaker. The wage stays elevated despite increased labor supply, sustaining the amplification.

\paragraph{Persistence and Propagation.}
The shock persists through two channels:

\textit{Exogenous persistence}: The AR(1) process for $A_t$ means high productivity today predicts high productivity tomorrow ($\rho_A > 0$).

\textit{Endogenous persistence}: Capital accumulation propagates the shock even after TFP returns to normal. High investment today raises the capital stock, which raises output tomorrow.

The strength of endogenous persistence depends on $\sigma$:
\begin{itemize}[nosep]
\item With $\sigma < 1$: A higher capital stock rapidly drives down the return to capital (strong diminishing returns). Investment falls quickly, limiting persistence.
\item With $\sigma > 1$: The return to capital falls slowly as capital accumulates. Investment stays elevated longer, increasing persistence.
\end{itemize}

\begin{calloutimportant}[Substitutes Generate More Persistence]
When labor and capital are substitutes ($\sigma > 1$), the RBC model generates more endogenous persistence. The return to capital remains elevated after a positive TFP shock, sustaining investment and prolonging the expansion. This helps match the observed persistence of business cycles without requiring extremely persistent exogenous shocks.
\end{calloutimportant}

\paragraph{Impulse Response Functions.}
The following table summarizes the qualitative impulse responses to a positive TFP shock:

\begin{center}
\renewcommand{\arraystretch}{1.3}
\begin{tabular}{@{}lccc@{}}
\toprule
\textbf{Variable} & \textbf{Impact} & \textbf{Persistence} ($\sigma < 1$) & \textbf{Persistence} ($\sigma > 1$) \\
\midrule
Output $y$ & $\uparrow$ & Returns quickly & Returns slowly \\
Consumption $c$ & $\uparrow$ & Gradual adjustment & Gradual adjustment \\
Investment $i$ & $\uparrow\uparrow$ & Falls quickly & Falls slowly \\
Labor $\ell$ & $\uparrow$ & Returns quickly & Returns slowly \\
Return $R$ & $\uparrow$ & Falls quickly & Falls slowly \\
Wage $w$ & $\uparrow$ & Returns quickly & Returns slowly \\
\bottomrule
\end{tabular}
\end{center}

% --------------------------------------------------------------
\subsection{Factor Shares Over the Business Cycle}
\label{subsec:rbc-shares}
% --------------------------------------------------------------

A key prediction of the CES framework concerns the cyclicality of factor shares. With Cobb-Douglas ($\sigma = 1$), factor shares are constant. With $\sigma \neq 1$, factor shares vary with the cycle.

\paragraph{The Labor Share Formula.}
From \cref{eq:cost-shares}, the labor share is:
\begin{equation}
s_L = \omega_L \left( \frac{w}{\Cstar} \right)^{1-\sigma} = \frac{w \cdot \ell}{y}.
\label{eq:labor-share-rbc}
\end{equation}

Using the marginal product formula \cref{eq:wage-ces}:
\begin{equation}
s_L = \omega_L^{1-\rho} A^\rho \left( \frac{y}{\ell} \right)^{1/\sigma - 1} = \omega_L^{1-\rho} A^\rho \left( \frac{y}{\ell} \right)^{(1-\sigma)/\sigma}.
\label{eq:labor-share-formula}
\end{equation}

\paragraph{Cyclicality of the Labor Share.}
Consider a positive TFP shock that raises output per worker $y/\ell$. The effect on the labor share depends on $(1-\sigma)/\sigma$:
\begin{itemize}[nosep]
\item If $\sigma < 1$: $(1-\sigma)/\sigma > 0$, so $s_L$ rises when $y/\ell$ rises. The labor share is \emph{procyclical}.
\item If $\sigma > 1$: $(1-\sigma)/\sigma < 0$, so $s_L$ falls when $y/\ell$ rises. The labor share is \emph{countercyclical}.
\item If $\sigma = 1$: $s_L = \omega_L$ regardless of $y/\ell$. The labor share is \emph{acyclical}.
\end{itemize}

\begin{calloutnote}[Empirical Evidence on Labor Share Cyclicality]
The cyclicality of the labor share is a subject of ongoing empirical debate. Early studies found a countercyclical labor share, suggesting $\sigma > 1$. More recent work, accounting for overhead labor and composition effects, finds evidence closer to acyclicality or mild procyclicality, suggesting $\sigma \leq 1$.

The answer may depend on the type of shock: demand shocks and supply shocks can have different implications for factor shares even with the same $\sigma$.
\end{calloutnote}

\paragraph{The Capital Share and Investment.}
The capital share $s_K = 1 - s_L$ inherits the opposite cyclicality. With $\sigma > 1$, the capital share is procyclical: booms are good times for capital owners.

This has implications for investment dynamics. When $\sigma > 1$, a positive shock raises the capital share, increasing the return to investment. This amplifies the investment response and contributes to the persistence discussed above.

% --------------------------------------------------------------
\subsection{Labor-Capital Substitution: The Automation Channel}
\label{subsec:rbc-automation}
% --------------------------------------------------------------

The standard RBC model features Hicks-neutral productivity shocks. We now consider the case where capital and labor are substitutes ($\sigma > 1$) and explore shocks that differentially affect factors---the ``automation'' channel.

\paragraph{Capital-Biased Technical Change.}
Suppose technological progress is embodied in capital: new capital goods are more productive than old ones. We model this as a fall in the effective price of capital (equivalently, an increase in capital-augmenting productivity $A_K$).

In the augmented production function:
\begin{equation}
y_t = \left( \omega_K^{1-\rho} (A_{K,t} \, k_t)^\rho + \omega_L^{1-\rho} \ell_t^\rho \right)^{1/\rho},
\label{eq:capital-biased}
\end{equation}
an increase in $A_{K,t}$ raises the effective capital stock.

\paragraph{Effects with $\sigma > 1$.}
When capital and labor are substitutes, capital-biased technical change has distinctive effects:

\begin{enumerate}[nosep]
\item \textbf{Output rises}: More effective capital increases production.
\item \textbf{The capital share rises}: With $\sigma > 1$, the factor receiving the productivity boost gains share.
\item \textbf{The labor share falls}: Workers receive a smaller fraction of a larger pie.
\item \textbf{The wage effect is ambiguous}: The wage $w = \partial y / \partial \ell$ may rise (due to higher output) or fall (due to substitution toward capital).
\end{enumerate}

\paragraph{The Automation Interpretation.}
If we interpret capital-biased technical change as ``automation''---machines becoming capable of performing tasks previously done by workers---then $\sigma > 1$ is the relevant case. Automation replaces labor, and the labor share declines.

The RBC framework with $\sigma > 1$ thus provides a foundation for analyzing automation-driven business cycles. Shocks to automation technology generate:
\begin{itemize}[nosep]
\item Expansions in output (good for aggregate welfare)
\item Declines in the labor share (potentially concerning for distribution)
\item Procyclical capital shares (booms favor capital owners)
\end{itemize}

\begin{calloutimportant}[RBC with Automation]
The standard RBC model with $\sigma = 1$ treats labor and capital symmetrically. Allowing $\sigma > 1$ introduces an asymmetry: capital can substitute for labor, and capital-biased shocks redistribute income toward capital owners.

This extension connects the RBC tradition to concerns about automation, inequality, and the declining labor share observed in many countries since the 1980s.
\end{calloutimportant}

% --------------------------------------------------------------
\subsection{Calibration and the CES Production Function}
\label{subsec:rbc-calibration}
% --------------------------------------------------------------

Quantitative RBC analysis requires calibrating the model parameters. The CES specification introduces $\sigma$ as a key parameter, in addition to the standard parameters $(\beta, \delta, \theta, \nu)$.

\paragraph{Standard Calibration Targets.}
The following moments are typically used to calibrate RBC models:

\begin{center}
\renewcommand{\arraystretch}{1.3}
\begin{tabular}{@{}lll@{}}
\toprule
\textbf{Parameter} & \textbf{Target} & \textbf{Typical Value} \\
\midrule
$\beta$ & Risk-free rate (4\% annual) & $0.99$ (quarterly) \\
$\delta$ & Investment-capital ratio & $0.025$ (quarterly) \\
$\theta$ & Consumption smoothing & $0.5$--$2$ \\
$\nu$ & Labor supply elasticity & $0.5$--$2$ \\
$\omega_K$ & Capital share & $0.33$ \\
\bottomrule
\end{tabular}
\end{center}

\paragraph{Calibrating $\sigma$.}
The elasticity of substitution $\sigma$ is harder to pin down. Identification comes from:

\textit{Long-run factor shares}: Trends in the labor share in response to changes in the capital-labor ratio or relative factor prices.

\textit{Cross-sectional variation}: Differences in factor shares across industries, countries, or time periods with different factor price ratios.

\textit{Business cycle moments}: The cyclicality of factor shares, the volatility of investment relative to output, and the persistence of output fluctuations.

Estimates vary widely:
\begin{itemize}[nosep]
\item Macro estimates (aggregate production functions): $\sigma \in [0.4, 1.0]$
\item Micro estimates (firm or industry level): $\sigma \in [0.5, 1.5]$
\item Recent estimates emphasizing capital heterogeneity: $\sigma > 1$ for equipment, $\sigma < 1$ for structures
\end{itemize}

\paragraph{The Nested CES Approach.}
Given the heterogeneity in estimates, the nested CES framework from \cref{sec:nested-capital} offers a resolution. Different capital types may have different substitution elasticities with labor:
\begin{equation}
y = A \cdot \M\left( \M(k_E, k_S; \bom^K, \sigma_K), \, \ell; \, \bom, \sigma \right).
\label{eq:nested-rbc}
\end{equation}

This allows, for example, $\sigma_K > 1$ (equipment and structures are substitutes) while $\sigma < 1$ (the capital composite and labor are complements). Such a specification can match both the apparent substitutability within capital and the complementarity between capital and labor at the aggregate level.

% --------------------------------------------------------------
\subsection{The RBC Model as a CES Economy}
\label{subsec:rbc-ces-summary}
% --------------------------------------------------------------

We conclude by summarizing the RBC model through the lens of CES aggregation.

\paragraph{The Three Margins.}
The RBC model involves optimization over three margins, each governed by a CES-type condition:

\begin{enumerate}[nosep]
\item \textbf{Intertemporal (Euler equation)}: Allocation of consumption across time, governed by the IES $\theta$.
\item \textbf{Intratemporal (labor supply)}: Allocation of time between work and leisure, governed by the Frisch elasticity $\nu$.
\item \textbf{Factors (production)}: Combination of capital and labor, governed by the elasticity of substitution $\sigma$.
\end{enumerate}

\paragraph{The Price Index Interpretation.}
Each margin has an associated ``price index'':
\begin{itemize}[nosep]
\item Intertemporal: The effective discount rate $\beta R$ determines the relative price of future consumption.
\item Intratemporal: The wage $w$ is the price of leisure in terms of consumption.
\item Factors: The unit cost $\Cstar$ aggregates factor prices.
\end{itemize}

\paragraph{Induced Weights and Shares.}
Optimal behavior generates shares as induced weights:
\begin{itemize}[nosep]
\item The saving rate is the ``share'' of income allocated to future consumption.
\item The labor share of time use is determined by the wage relative to the marginal utility of leisure.
\item Factor shares are the cost shares $s_K$ and $s_L$.
\end{itemize}

\begin{table}[H]
\centering
\caption{The RBC Model: Key Equations and CES Structure}
\label{tab:rbc-summary}
\renewcommand{\arraystretch}{1.6}
\begin{tabular}{@{}ll@{}}
\toprule
\textbf{Component} & \textbf{Formula} \\
\midrule
Production & $y = A \left( \omega_K^{1-\rho} k^\rho + \omega_L^{1-\rho} \ell^\rho \right)^{1/\rho}$ \\[6pt]
Euler equation & $c_t^{-1/\theta} = \beta \, \E_t \left[ c_{t+1}^{-1/\theta} R_{t+1} \right]$ \\[6pt]
Labor supply & $\chi \ell^{1/\nu} c^{1/\theta} = w$ \\[6pt]
Return to capital & $R = \omega_K^{1-\rho} A^\rho (y/k)^{1/\sigma} + 1 - \delta$ \\[6pt]
Wage & $w = \omega_L^{1-\rho} A^\rho (y/\ell)^{1/\sigma}$ \\[6pt]
Labor share & $s_L = \omega_L (w/\Cstar)^{1-\sigma}$ \\[6pt]
\bottomrule
\end{tabular}
\end{table}

\begin{callouttip}[Key Insight]
The RBC model is a dynamic CES economy. Productivity shocks propagate through the equilibrium conditions, with the elasticity of substitution $\sigma$ governing:
\begin{itemize}[nosep]
\item \textbf{Amplification}: How strongly labor responds to wage changes
\item \textbf{Persistence}: How slowly returns fall as capital accumulates
\item \textbf{Distribution}: How factor shares move over the cycle
\end{itemize}
The Cobb-Douglas case ($\sigma = 1$) is a knife-edge: factor shares are constant, and the model has specific persistence properties. Departing from $\sigma = 1$ enriches the model's predictions and connects to debates about automation, inequality, and the declining labor share.

The nested CES extension allows for heterogeneous capital, providing a bridge to the richer models developed in the next section.
\end{callouttip}


% ==============================================================
\section{Beyond RBC: Heterogeneity}
\label{sec:beyond-rbc}
% ==============================================================

The representative agent RBC model provides a powerful framework for understanding aggregate fluctuations, but it abstracts from the rich heterogeneity observed in data. Workers differ in skills, capital goods differ in their characteristics, and these differences interact in economically important ways. This section extends the CES framework to incorporate heterogeneity in both labor and capital, connecting to the discrete choice and recursive structures developed in earlier chapters.

% --------------------------------------------------------------
\subsection{Why Heterogeneity Matters}
\label{subsec:why-heterogeneity}
% --------------------------------------------------------------

The representative agent assumption imposes strong restrictions. All workers are identical, all capital goods are identical, and shocks affect everyone symmetrically. Relaxing these assumptions opens new channels for understanding macroeconomic dynamics.

\paragraph{Heterogeneous Labor.}
Workers differ along many dimensions: education, experience, occupation, and innate ability. These differences matter for:
\begin{itemize}[nosep]
\item \textbf{Wage inequality}: Different skill types command different wages.
\item \textbf{Skill premia}: The college wage premium has risen dramatically since 1980.
\item \textbf{Differential exposure to shocks}: Recessions disproportionately affect less-skilled workers.
\item \textbf{Skill-biased technical change}: Technology may complement some skills while substituting for others.
\end{itemize}

\paragraph{Heterogeneous Capital.}
Capital goods also differ: structures versus equipment, computers versus vehicles, tangible versus intangible. These differences matter for:
\begin{itemize}[nosep]
\item \textbf{Investment composition}: The share of equipment in total investment has risen over time.
\item \textbf{Differential depreciation}: Equipment depreciates faster than structures.
\item \textbf{Sector-specific technical change}: The price of computing equipment has fallen dramatically.
\item \textbf{Complementarity patterns}: Some capital types may complement skilled labor while substituting for unskilled labor.
\end{itemize}

\paragraph{The CES Approach to Heterogeneity.}
The CES framework provides a natural language for modeling heterogeneity. Different types of labor or capital are aggregated using generalized means, with the elasticity of substitution governing how easily one type can replace another. The nested CES structure allows for rich patterns of complementarity and substitutability across different levels of aggregation.

% --------------------------------------------------------------
\subsection{Heterogeneous Labor: Skill Types}
\label{subsec:heterogeneous-labor}
% --------------------------------------------------------------

We begin by introducing multiple skill types into the production function. This extension captures the observation that workers with different skills are imperfect substitutes.

\paragraph{The Labor Aggregate.}
Suppose there are two types of labor: skilled ($L_S$) and unskilled ($L_U$). The labor aggregate is a CES composite:
\begin{equation}
L = \M(L_S, L_U; \bom^L, \sigma_L) = \left( \omega_S^{1-\rho_L} L_S^{\rho_L} + \omega_U^{1-\rho_L} L_U^{\rho_L} \right)^{1/\rho_L},
\label{eq:labor-aggregate}
\end{equation}
where $\sigma_L = 1/(1-\rho_L)$ is the elasticity of substitution between skill types, and $\omega_S + \omega_U = 1$.

\paragraph{The Production Function.}
Output combines the labor aggregate with capital:
\begin{equation}
Y = A \cdot \M(K, L; \bom, \sigma) = A \left( \omega_K^{1-\rho} K^\rho + \omega_L^{1-\rho} L^\rho \right)^{1/\rho}.
\label{eq:production-hetero-labor}
\end{equation}

This is a nested CES structure with labor types in the inner nest and capital-labor in the outer nest.

\paragraph{Wages by Skill Type.}
The wage of each skill type is its marginal product:
\begin{align}
w_S &= \frac{\partial Y}{\partial L_S} = \frac{\partial Y}{\partial L} \cdot \frac{\partial L}{\partial L_S}, \label{eq:wage-skilled} \\[6pt]
w_U &= \frac{\partial Y}{\partial L_U} = \frac{\partial Y}{\partial L} \cdot \frac{\partial L}{\partial L_U}. \label{eq:wage-unskilled}
\end{align}

Using the CES formulas:
\begin{equation}
w_S = \omega_L^{1-\rho} A^\rho \left( \frac{Y}{L} \right)^{1/\sigma} \cdot \omega_S^{1-\rho_L} \left( \frac{L}{L_S} \right)^{1/\sigma_L}.
\label{eq:wage-skilled-full}
\end{equation}

\paragraph{The Skill Premium.}
The ratio of skilled to unskilled wages is the skill premium:
\begin{equation}
\frac{w_S}{w_U} = \frac{\omega_S}{\omega_U} \left( \frac{L_S}{L_U} \right)^{-1/\sigma_L} = \frac{\omega_S}{\omega_U} \left( \frac{L_U}{L_S} \right)^{1/\sigma_L}.
\label{eq:skill-premium}
\end{equation}

The skill premium depends on:
\begin{itemize}[nosep]
\item Relative weights $\omega_S/\omega_U$: reflecting productivity differences
\item Relative supplies $L_S/L_U$: more skilled workers reduce the premium
\item The elasticity $\sigma_L$: lower elasticity means supplies matter more
\end{itemize}

\begin{calloutimportant}[The Skill Premium and $\sigma_L$]
If $\sigma_L < \infty$, skilled and unskilled workers are imperfect substitutes. An increase in skilled labor supply reduces the skill premium, but the effect is larger when $\sigma_L$ is smaller (less substitutable).

Estimates of $\sigma_L$ typically range from 1.4 to 2.0, implying significant but not perfect substitutability. This means that supply shifts (e.g., increased college enrollment) have substantial effects on wage inequality.
\end{calloutimportant}

\paragraph{Skill-Biased Technical Change.}
Suppose technological progress augments skilled labor: $L_S$ is replaced by $A_S \cdot L_S$ in the production function. The skill premium becomes:
\begin{equation}
\frac{w_S}{w_U} = \frac{\omega_S}{\omega_U} \left( \frac{A_S L_S}{L_U} \right)^{(\sigma_L - 1)/\sigma_L} \cdot A_S.
\label{eq:skill-premium-sbtc}
\end{equation}

When $\sigma_L > 1$, an increase in $A_S$ raises the skill premium through two channels:
\begin{enumerate}[nosep]
\item Direct effect: skilled workers are more productive (factor $A_S$)
\item Indirect effect: effective skilled labor rises, but with $\sigma_L > 1$, this raises relative demand for skilled workers
\end{enumerate}

This is the canonical model of skill-biased technical change (SBTC), used to explain the rising college wage premium since 1980.

% --------------------------------------------------------------
\subsection{Heterogeneous Capital: Equipment and Structures}
\label{subsec:heterogeneous-capital}
% --------------------------------------------------------------

We now extend the model to include multiple capital types, building on the nested CES framework from \cref{sec:nested-capital}.

\paragraph{The Capital Aggregate.}
Capital consists of equipment ($K_E$) and structures ($K_S$):
\begin{equation}
K = \M(K_E, K_S; \bom^K, \sigma_K) = \left( \omega_E^{1-\rho_K} K_E^{\rho_K} + \omega_S^{1-\rho_K} K_S^{\rho_K} \right)^{1/\rho_K}.
\label{eq:capital-aggregate}
\end{equation}

\paragraph{The Full Nested Structure.}
Combining heterogeneous labor and heterogeneous capital:
\begin{equation}
Y = A \cdot \M\left( \M(K_E, K_S; \bom^K, \sigma_K), \, \M(L_S, L_U; \bom^L, \sigma_L); \, \bom, \sigma \right).
\label{eq:full-nested}
\end{equation}

This three-level nested CES has:
\begin{itemize}[nosep]
\item Inner capital nest: elasticity $\sigma_K$ between equipment and structures
\item Inner labor nest: elasticity $\sigma_L$ between skilled and unskilled labor
\item Outer nest: elasticity $\sigma$ between the capital and labor composites
\end{itemize}

\paragraph{Price Indices at Each Level.}
The nested structure generates nested price indices:
\begin{align}
C_K^* &= \M_{1-\sigma_K}(r_E, r_S; \bom^K), \label{eq:capital-price-index} \\[6pt]
W^* &= \M_{1-\sigma_L}(w_S, w_U; \bom^L), \label{eq:labor-price-index} \\[6pt]
\Cstar &= \M_{1-\sigma}(C_K^*, W^*; \bom). \label{eq:total-price-index}
\end{align}

The labor price index $W^*$ is the unit cost of the labor composite---a CES aggregate of wages.

% --------------------------------------------------------------
\subsection{Capital-Skill Complementarity}
\label{subsec:capital-skill-complementarity}
% --------------------------------------------------------------

A key hypothesis in the inequality literature is that capital (especially equipment) is more complementary with skilled labor than with unskilled labor. The nested CES framework can capture this through alternative nesting structures.

\paragraph{The Krusell-Ohanian-Rios-Rull-Violante (KORV) Specification.}
A leading specification places equipment and unskilled labor in one nest, and this composite with skilled labor in another:
\begin{equation}
Y = A \cdot \M\left( K_S, \, \M\left( L_S, \, \M(K_E, L_U; \bom^{EU}, \sigma_{EU}); \bom^{S}, \sigma_S \right); \bom, \sigma_{\text{out}} \right).
\label{eq:korv}
\end{equation}

The key feature is that equipment ($K_E$) and unskilled labor ($L_U$) are in the innermost nest with elasticity $\sigma_{EU}$. If $\sigma_{EU} > 1$, equipment substitutes for unskilled labor.

\paragraph{The Capital-Skill Complementarity Hypothesis.}
The KORV specification implies:
\begin{itemize}[nosep]
\item Equipment and unskilled labor are substitutes ($\sigma_{EU} > 1$)
\item The equipment-unskilled composite and skilled labor are complements ($\sigma_S < 1$)
\item Cheaper equipment replaces unskilled workers and raises demand for skilled workers
\end{itemize}

This generates a link between the falling price of equipment and rising wage inequality.

\begin{calloutnote}[Automation and Inequality]
The capital-skill complementarity framework formalizes the intuition that automation displaces some workers while complementing others. As equipment becomes cheaper:
\begin{enumerate}[nosep]
\item Firms substitute equipment for unskilled labor (inner nest, $\sigma_{EU} > 1$)
\item The effective supply of the equipment-unskilled composite rises
\item This raises demand for skilled labor (outer nest, $\sigma_S < 1$)
\item The skill premium rises
\end{enumerate}
This mechanism can explain both the falling labor share (especially for unskilled workers) and the rising skill premium observed since 1980.
\end{calloutnote}

\paragraph{Quantitative Implications.}
Krusell et al. (2000) estimate the KORV model and find:
\begin{itemize}[nosep]
\item $\sigma_{EU} \approx 1.67$: equipment and unskilled labor are substitutes
\item $\sigma_S \approx 0.67$: the composite and skilled labor are complements
\item The falling relative price of equipment can explain much of the rise in the skill premium
\end{itemize}

% --------------------------------------------------------------
\subsection{Connecting to Discrete Choice}
\label{subsec:discrete-choice-connection}
% --------------------------------------------------------------

The CES framework with heterogeneous factors connects deeply to discrete choice models. This connection, developed in \cref{sec:discrete-choice}, provides both economic intuition and computational tools.

\paragraph{The CES-Logit Isomorphism.}
Recall that the CES demand system is isomorphic to the multinomial logit. In the production context:
\begin{itemize}[nosep]
\item Factor shares $s_i$ correspond to choice probabilities $\pi_i$
\item The elasticity $\sigma$ corresponds to the inverse scale parameter $1/\lambda$
\item Factor prices $p_i$ correspond to (negative) utilities $-v_i$
\end{itemize}

The CES share formula:
\begin{equation}
s_i = \frac{\omega_i \, p_i^{1-\sigma}}{\sum_j \omega_j \, p_j^{1-\sigma}}
\label{eq:ces-share}
\end{equation}
is equivalent to the logit probability:
\begin{equation}
\pi_i = \frac{\exp(v_i / \lambda)}{\sum_j \exp(v_j / \lambda)}
\label{eq:logit-prob}
\end{equation}
under the transformation $v_i = -\lambda(1-\sigma) \log p_i + \lambda \log \omega_i$.

\paragraph{Nested Logit and Nested CES.}
The nested CES structure corresponds to the nested logit model:
\begin{itemize}[nosep]
\item Groups of factors (e.g., skill types within labor) correspond to nests
\item Different elasticities at different levels correspond to different scale parameters
\item The inclusive value (log-sum) in nested logit corresponds to the price index in nested CES
\end{itemize}

\begin{calloutimportant}[The Discrete Choice Interpretation]
Think of the firm as ``choosing'' how to allocate expenditure across factors. The CES structure implies that this choice follows a logit model:
\begin{itemize}[nosep]
\item Cheaper factors are more likely to be ``chosen'' (higher expenditure share)
\item The elasticity $\sigma$ governs how sensitive the choice is to price
\item Nested CES allows for correlation within groups (skilled vs. unskilled labor)
\end{itemize}
This interpretation is useful for building intuition and for computational methods (simulated method of moments, BLP-style estimation).
\end{calloutimportant}

\paragraph{Random Coefficients and Continuous Heterogeneity.}
The discrete choice perspective suggests extensions to continuous heterogeneity. Instead of a finite number of skill types, workers may be distributed along a continuum of skills. The CES aggregate becomes an integral:
\begin{equation}
L = \left( \int_0^1 \omega(z)^{1-\rho_L} L(z)^{\rho_L} \, dz \right)^{1/\rho_L},
\label{eq:continuous-ces}
\end{equation}
where $z$ indexes skill types.

This connects to the ``task-based'' approach to labor markets, where workers of different skills have comparative advantage in different tasks, and automation changes which tasks are performed by machines versus humans.

% --------------------------------------------------------------
\subsection{Connecting to Recursive Structures}
\label{subsec:recursive-connection}
% --------------------------------------------------------------

The dynamic models with heterogeneity connect to the recursive structures developed in \cref{sec:recursive-decomposition}. The value function inherits the CES structure, enabling tractable characterization of equilibrium.

\paragraph{The Bellman Equation with Heterogeneity.}
Consider a planner's problem with heterogeneous capital:
\begin{equation}
V(k_E, k_S, A) = \max_{c, i_E, i_S} \left\{ u(c) + \beta \, \E \left[ V(k_E', k_S', A') \right] \right\}
\label{eq:bellman-hetero}
\end{equation}
subject to:
\begin{align}
c + i_E + i_S &= Y(k_E, k_S, A), \\
k_E' &= (1-\delta_E) k_E + i_E, \\
k_S' &= (1-\delta_S) k_S + i_S.
\end{align}

\paragraph{Homogeneity and Dimension Reduction.}
With CES production and CRRA utility, the value function is homogeneous:
\begin{equation}
V(k_E, k_S, A) = \phi(k_E/k_S, A) \cdot \frac{(k_E + k_S)^{1-1/\theta}}{1 - 1/\theta}.
\label{eq:value-homogeneous}
\end{equation}

This reduces the state space: instead of tracking $(k_E, k_S)$ separately, we can track total capital and the composition ratio $k_E/k_S$. The CES structure ensures that the composition dynamics are tractable.

\paragraph{The Recursive CES Structure.}
The nested CES production function has a recursive structure that mirrors the recursive structure of the value function. Define:
\begin{align}
K &= \M(k_E, k_S; \bom^K, \sigma_K), \\
Y &= A \cdot \M(K, L; \bom, \sigma).
\end{align}

The first-order conditions for investment allocation satisfy:
\begin{equation}
\frac{\partial V / \partial k_E}{\partial V / \partial k_S} = \frac{1 - \delta_E}{1 - \delta_S} \cdot \frac{\partial Y / \partial k_E}{\partial Y / \partial k_S}.
\label{eq:investment-foc}
\end{equation}

With CES, this ratio depends only on the composition $k_E/k_S$ and the depreciation rates, not on the level of capital. This separability is a hallmark of the CES structure.

\paragraph{Epstein-Zin with Heterogeneous Capital.}
The Epstein-Zin framework from \cref{sec:epstein-zin} extends naturally to heterogeneous capital. The value function becomes:
\begin{equation}
V_t = \M\left( c_t, \, \CE\left( V_{t+1}(k_E', k_S', A') \right); \bom^\beta, \theta \right),
\label{eq:ez-hetero}
\end{equation}
where $\CE$ is the certainty equivalent over future states.

This ``CES over CES'' structure---CES aggregation over time, CES aggregation over capital types, CES certainty equivalent over states---maintains tractability while allowing rich dynamics.

% --------------------------------------------------------------
\subsection{Background Risk and Precautionary Behavior}
\label{subsec:background-risk}
% --------------------------------------------------------------

Heterogeneity interacts with risk in important ways. When agents face uninsurable ``background risk,'' their behavior changes even with respect to insurable risks.

\paragraph{Background Risk in the Labor Market.}
Workers face idiosyncratic income risk that cannot be fully insured:
\begin{equation}
y_t^i = w_t \cdot \ell_t^i \cdot \exp(\eta_t^i), \qquad \eta_t^i \sim N(-\sigma_\eta^2/2, \sigma_\eta^2),
\label{eq:idiosyncratic-income}
\end{equation}
where $\eta_t^i$ is an idiosyncratic shock to worker $i$'s effective labor supply.

\paragraph{Precautionary Saving.}
Background risk induces precautionary saving. With CRRA utility, the Euler equation becomes:
\begin{equation}
c_t^{-1/\theta} = \beta R \, \E_t \left[ c_{t+1}^{-1/\theta} \right] > \beta R \, \left( \E_t[c_{t+1}] \right)^{-1/\theta},
\label{eq:precautionary-euler}
\end{equation}
where the inequality follows from Jensen's inequality when $1/\theta > 0$ (risk aversion).

The gap represents precautionary saving: agents save more than they would without risk, to self-insure against future income shocks.

\paragraph{Effects on Interest Rates.}
In general equilibrium, precautionary saving lowers the interest rate. The steady-state return satisfies:
\begin{equation}
R^* = \frac{(1+g)^{1/\theta}}{\beta} \cdot \underbrace{\frac{1}{\E[m^{-1/\theta}] / (\E[m])^{-1/\theta}}}_{\text{precautionary factor} < 1},
\label{eq:interest-rate-precautionary}
\end{equation}
where $m$ captures the distribution of marginal utilities.

Greater background risk (higher $\sigma_\eta$) reduces $R^*$: households want to save more, driving down the equilibrium return.

\begin{calloutnote}[The Risk-Free Rate Puzzle]
The low level of risk-free interest rates has been a puzzle in asset pricing. One resolution involves background risk: if households face substantial uninsurable income risk, they have strong precautionary motives, which can explain low equilibrium interest rates even with moderate time preference.

The CES framework with heterogeneity provides a tractable setting to study these effects quantitatively.
\end{calloutnote}

% --------------------------------------------------------------
\subsection{Shocks with Heterogeneous Factors}
\label{subsec:shocks-heterogeneity}
% --------------------------------------------------------------

We conclude by examining how aggregate shocks propagate when factors are heterogeneous. The key insight is that different shocks have different distributional implications.

\paragraph{Neutral vs. Biased Shocks.}
Consider three types of productivity shocks:
\begin{enumerate}[nosep]
\item \textbf{Hicks-neutral}: $A$ rises, scaling output proportionally
\item \textbf{Capital-biased}: $A_K$ rises, augmenting effective capital
\item \textbf{Skill-biased}: $A_S$ rises, augmenting effective skilled labor
\end{enumerate}

Each has different implications for factor shares and inequality.

\paragraph{Impulse Responses by Factor Type.}
The following table summarizes the qualitative responses to each shock type:

\begin{center}
\renewcommand{\arraystretch}{1.3}
\begin{tabular}{@{}lccc@{}}
\toprule
& \textbf{Hicks-Neutral} & \textbf{Capital-Biased} & \textbf{Skill-Biased} \\
\midrule
Output $Y$ & $\uparrow$ & $\uparrow$ & $\uparrow$ \\
Capital share $s_K$ & $\approx$ constant & $\uparrow$ (if $\sigma > 1$) & $\downarrow$ \\
Labor share $s_L$ & $\approx$ constant & $\downarrow$ (if $\sigma > 1$) & $\uparrow$ \\
Skill premium $w_S/w_U$ & $\approx$ constant & $\uparrow$ (if CSC) & $\uparrow$ (if $\sigma_L > 1$) \\
\bottomrule
\end{tabular}
\end{center}

(CSC = capital-skill complementarity)

\paragraph{Business Cycles with Distributional Consequences.}
In the heterogeneous-factor model, business cycles have distributional consequences beyond the representative agent framework:
\begin{itemize}[nosep]
\item Recessions may disproportionately affect unskilled workers (if driven by capital-biased shocks)
\item Booms may increase inequality (if driven by skill-biased shocks)
\item The welfare cost of fluctuations differs across worker types
\end{itemize}

\paragraph{Policy Implications.}
Understanding the source of shocks matters for policy:
\begin{itemize}[nosep]
\item Neutral shocks: stabilization policy affects all factors symmetrically
\item Biased shocks: stabilization policy has distributional implications
\item Targeted policies (e.g., training programs, investment incentives) can address factor-specific disruptions
\end{itemize}

% --------------------------------------------------------------
\subsection{Summary}
\label{subsec:beyond-rbc-summary}
% --------------------------------------------------------------

\begin{table}[H]
\centering
\caption{Beyond RBC: Heterogeneity Extensions}
\label{tab:beyond-rbc-summary}
\renewcommand{\arraystretch}{1.6}
\begin{tabular}{@{}ll@{}}
\toprule
\textbf{Extension} & \textbf{Key Feature} \\
\midrule
Heterogeneous labor & $L = \M(L_S, L_U; \sigma_L)$, skill premium \\[6pt]
Heterogeneous capital & $K = \M(K_E, K_S; \sigma_K)$, investment composition \\[6pt]
Capital-skill complementarity & Equipment substitutes for unskilled, complements skilled \\[6pt]
Discrete choice connection & CES shares $\leftrightarrow$ logit probabilities \\[6pt]
Recursive structure & Homogeneity reduces state space \\[6pt]
Background risk & Precautionary saving, lower interest rates \\[6pt]
\bottomrule
\end{tabular}
\end{table}

\begin{callouttip}[Key Insight]
Moving beyond the representative agent RBC model, the CES framework accommodates rich heterogeneity:
\begin{itemize}[nosep]
\item \textbf{Multiple factor types}: Skilled/unskilled labor, equipment/structures capital
\item \textbf{Nested substitution patterns}: Different elasticities at different levels
\item \textbf{Capital-skill complementarity}: A unified explanation for falling labor shares and rising skill premia
\item \textbf{Connections}: To discrete choice (CES-logit isomorphism) and recursive methods (homogeneity)
\end{itemize}

The nested CES structure is more than a tractable functional form---it encodes economically meaningful hypotheses about how different factors interact. Estimating the elasticities at each level of the nest is an active area of empirical research, with implications for understanding automation, inequality, and the future of work.
\end{callouttip}


